% =====================================================================
%  content.tex --- body of the notes
%  Translated from the Polish original. Substantive issues carried over
%  from the source are marked with "% NOTE:" comments. Those that are still
%  WRONG OR DOUBTFUL AS PRINTED additionally carry a \merr marker, which
%  renders as a red (!) in the PDF. NOTEs without a marker record a source
%  typo that has already been repaired here.
% =====================================================================

\section{Inertial frames, principles of relativity, the Galilean transformation, the Lorentz transformation and its consequences}

An inertial frame is a frame of reference with respect to which every body, not subject to external interactions with other bodies, moves in uniform rectilinear motion or remains at rest. The existence of such a frame is postulated by Newton's first law of motion. An inertial frame of reference may equivalently be defined as a frame in which no fictitious inertial forces appear.

The Galilean principle of relativity states that the laws of motion are identical in all inertial frames of reference, i.e.\ there is no privileged frame of reference. It was formulated in 1604 in the form of the following law: all frames of reference moving with respect to one another at constant velocity (including $\vec{v}=0$) are equivalent. Its consequence is the set of Galilean transformations, in which time and the distance between any two points remain constant regardless of the frame of reference.

Einstein additionally postulated that in every inertial frame of reference the speed of light must have the same value. He based this assumption on an analysis of the transformations that leave Maxwell's equations invariant. The Galilean transformations were thereby generalised to the Lorentz transformation to accommodate the Einstein's postulate.

Consider two frames of reference: the first at rest and the second---the primed one---moving with velocity $v$ along the $x$ axis. The Lorentz transformation between these two frames takes the following form:

\begin{align*}
    t' &= \frac{t-xv/c^2}{\sqrt{1-v^2/c^2}} \\
    x' &= \frac{x-tv}{\sqrt{1-v^2/c^2}} \\
    y' &= y \\
    z' &= z
\end{align*}

The following notation may also be adopted:

\begin{equation*}
    \gamma = \frac{1}{\sqrt{1 - \beta^2}},\; \beta=\frac{v}{c},
\end{equation*}

which yields:

\begin{align*}
    t' &= \gamma\cdot(t-xv/c^2) \\
    x' &= \gamma\cdot(x-tv) \\
    y' &= y \\
    z' &= z
\end{align*}

The parameter $\gamma$ is also known as the Lorentz factor, and strictly: $\gamma > 1$. Note that in the limit $c \xrightarrow{} \infty$ the transformation reduces to the Galilean transformation.

\needspace{4\baselineskip}
The consequences of the Lorentz transformation include:
\begin{enumerate}
    \item Time dilation, described by the formula $\Delta t = \gamma\Delta t_0$, where $\Delta t_0$ is the duration of a phenomenon recorded by an observer at rest with respect to that phenomenon, and $\Delta t$ is the duration of the same phenomenon recorded by another observer moving with velocity $v$ with respect to the first. The event appears \textit{shortest} observed in its own frame.
    \item Lorentz contraction, described by the formula $L = {L_0}/{\gamma}$, where $L_0$ is the length of an object at rest and $L$ is the length of the object recorded by an observer moving with respect to it at velocity $v$. The object appears \textit{longest} in in its own frame.
    \item Relativity of simultaneity: two events determined by one observer to be simultaneous need not be simultaneous for another observer. Simultaneity translates to other frames only for the events that happen at the same point in space.
    \item The description of vector quantities in four dimensions. Position, momentum and potential are described respectively by a four-vector, a four-momentum and a four-potential. For example:
    % Four-momentum as an explicit column vector
    \begin{equation*}
        p^{\mu} =
        \begin{bmatrix}
            p^{0} \\ p^{1} \\ p^{2} \\ p^{3}
        \end{bmatrix}
        =
        \begin{bmatrix}
            E/c \\ p_x \\ p_y \\ p_z
        \end{bmatrix}
    \end{equation*}
    
\end{enumerate}

\section{Fundamental interactions; carriers and range of interactions, charges and characteristic energy scales}

Four types of fundamental interaction are currently known, each mediated by its own carrier:

\begin{itemize}
    \item the strong interaction
    \item the weak interaction
    \item the electromagnetic interaction
    \item the gravitational interaction
\end{itemize}

In quantum theories, interactions are transmitted by certain particles called interaction carriers or interaction quanta. There are two categories of such particles---real particles (transmitting the fundamental interactions) and quasiparticles (pseudoparticles: the phonon, the plasmon, the magnon).

\begin{table}[h!]
\centering
\begin{tabular}{|c|c|c|c|c|c|}
\hline
interaction & carrier & range & charge & mass & spin \\
\hline
strong & gluon & $10^{-15}$\,m & colour & 0 & 1 \\
\hline
weak & $W^+, W^-, Z^0$ boson & $10^{-18}$\,m & isospin & $m_0>0$ & 1 \\
\hline
electromagnetic & photon & $\infty$ & --- & 0 & s=1 \\
\hline
gravitational & graviton (hypothetical) & $\infty$ & --- & 0 & s=2 \\
\hline
\end{tabular}
\end{table}

% NOTE: the source writes the uncertainty relation with "=" and omits the
% inequality; strictly $\Delta E \Delta t \geq \hbar/2$.
It follows from the Heisenberg uncertainty principle, $\Delta E \Delta t = \frac{\hbar}{2}$\merr, that field quanta may exist for a finite time set by the magnitude of these energies, such that the product of these quantities is of the order of the Planck constant. The uncertainty principle allows the range of an interaction to be estimated, this range depending on the rest mass of the carrier particle. The masslessness of photons, as well as the expected masslessness of gravitons, is on this account the reason why these forces are long-ranged. % NOTE: the source reads "particles of zero energy, and hence rest mass" --- it should read "of zero rest mass".
Particles of zero energy, and hence of zero rest mass, would have infinite range.\merr

Intermediate bosons are exchanged by electrons, neutrinos and other particles participating in the weak interaction during collisions. The $W$ boson occurs in two basic forms: the particle $W^+$ and its antiparticle $W^-$. Both have the same spin and mass, $80.385 \pm 0.015$~GeV, and differ only in electric charge. The $W$ bosons carry one of the charges of the electroweak interaction, namely weak isospin. Its value is $+1$ and $-1$ for the $W^+$ and $W^-$ bosons respectively. The most important phenomenon involving $W$ bosons is beta decay. In this process a down quark (d) contained in a neutron decays into an up quark (u) and a $W$ boson, which in turn decays into a lepton--neutrino pair for a given lepton (more precisely: $W^-$ decays into a lepton plus an antineutrino, and $W^+$ into an antilepton plus a neutrino). As the energy increases, these leptons may be the electron, the muon or the tau. The $Z$ boson is simultaneously its own antiparticle. The mass (rest energy) of the $Z^0$ boson is $91.1876 \pm 0.0021$~GeV, and it is electrically neutral. The $Z$ boson decays into pairs of charged leptons (electron--antielectron, muon--antimuon) and into pairs of quarks or neutrinos (neutrino--antineutrino).

The graviton is a hypothetical particle whose existence is assumed by the theory of quantum gravity based on quantum field theory, but not on the Standard Model. The graviton has neither mass nor electric charge, and is a boson of spin $s = 2$.

There is no physical law implying that the detection of a graviton is impossible, but the required detector mass makes it impossible in practice. For example, it is estimated that a detector of Jupiter's mass operating at 100\% efficiency would be capable of detecting one graviton in ten years. The reason for this is the very small cross-section of the graviton for interaction with matter. A further problem is shielding from background noise. The two basic sources of noise are cosmic radiation and neutrinos.

The gluon is the carrier of the strong interaction, which means that this interaction consists in the exchange of gluons between quarks or between other gluons. The theory describing the interactions of gluons and quarks is quantum chromodynamics. The gluon carries colour charge. Gluons exchanged between quarks form a field of colour forces. Gluons have their own colour, and therefore, besides transmitting interactions between quarks, they are able to interact with themselves, which leads for instance to the appearance of so-called gluon loops in Feynman diagrams.

% NOTE: the following paragraph is internally inconsistent (as in the source):
% a gluon carries a colour AND an anticolour; the claim that it carries
% "only one colour charge" is incorrect.
As carriers of the strong interaction, gluons have double charges: one colour and one anticolour. However, each gluon carries only one colour charge: either one colour or one anticolour.\merr For example, there exists a red--antiblue gluon (red/antiblue gluon), which carries the colour red. When a quark emits or absorbs a gluon, the colour of the quark changes. Suppose, for instance, that a red quark turns into a blue quark and emits a red--antiblue gluon. Such a gluon is absorbed by another blue quark. As a result of absorbing the red--antiblue gluon, the blue quark turns into a red quark.

\section{Conservation laws in physics and their connection with symmetry laws}

The symmetries of physical systems are closely connected with conservation laws. Noether's theorem formalises this connection, stating that to every continuous symmetry of the action there corresponds a conservation law.

\subsection{Noether's theorem}

Noether's theorem states that if the action \( S \) is invariant with respect to some group of continuous transformations, then there exists a physical quantity that is conserved. Mathematically, if

\begin{equation}
\delta S = 0,
\end{equation}

then there exists a conserved current \( j^\mu \) such that

\begin{equation}
\partial_\mu j^\mu = 0.
\end{equation}

\subsection{Examples of the application of Noether's theorem}

\subsubsection{Conservation of energy}

Time symmetry (homogeneity of time) leads to the law of conservation of energy. If the Lagrangian \( L \) does not depend explicitly on time, then energy is conserved:

\begin{equation}
H = \frac{\partial L}{\partial \dot{q}} \dot{q} - L = \text{const}.
\end{equation}

\subsubsection{Conservation of momentum}

Spatial symmetry (homogeneity of space) leads to the law of conservation of momentum. If the Lagrangian does not depend explicitly on position, then momentum is conserved:

\begin{equation}
\vec{p} = \frac{\partial L}{\partial \dot{\vec{q}}} = \text{const}.
\end{equation}

\subsubsection{Conservation of angular momentum}

Rotational symmetry (isotropy of space) leads to the law of conservation of angular momentum. If the Lagrangian is invariant with respect to rotations, then angular momentum is conserved:

\begin{equation}
\vec{L} = \vec{q} \times \vec{p} = \text{const}.
\end{equation}

\section{Newton's laws of motion and the limits of their applicability; classical mechanics as a limiting form of quantum mechanics; non-relativistic classical mechanics as the limit of relativistic classical mechanics}

Newton's laws of motion are three laws lying at the foundation of classical mechanics, specifying the relations between the motion of a body and the forces acting on it, which is why they are also called the laws of motion. In quantum mechanics they are for the most part inapplicable, and in relativistic mechanics they hold only within a limited range.

\subsection{Newton's first law}

In an inertial frame of reference, if no force acts on a body, or if the acting forces balance one another, then the body remains at rest or moves in uniform rectilinear motion.

An inertial frame is a frame of reference with respect to which every body not subject to external interactions with other bodies moves without acceleration (i.e.\ in uniform rectilinear motion) or remains at rest.

The first law assumes the existence of an inertial frame of reference but does not indicate how such a frame is to be sought, and it is therefore a postulate of the existence of an inertial frame of reference, and is formulated as follows:

\textit{There exists a frame of reference in which a particle not subject to interaction with its surroundings remains at rest or moves along a straight line with constant velocity.}

\subsection{Newton's second law}

If an unbalanced force acts on a body, then the body moves with an acceleration directly proportional to the force and inversely proportional to the mass of the body.

\[
\vec{F} = m \vec{a}
\]

where:
\begin{itemize}
    \item \(\vec{F}\) --- the force acting on the body,
    \item \(m\) --- the mass of the body,
    \item \(\vec{a}\) --- the acceleration of the body.
\end{itemize}

\subsection{Newton's third law}

Interactions between bodies are always mutual. In an inertial frame of reference the forces of mutual interaction of two bodies have the same magnitude, the same direction, opposite senses and different points of application (each acts on a different body).

\[
\vec{F}_{12} = -\vec{F}_{21}
\]

In classical mechanics, Newton's third law and the law of conservation of momentum are for the most part equivalent. Neither holds when the physical system is not isolated, or during the creation of virtual particle pairs. The laws of motion can also be written for angular quantities in rotational motion, but a straightforward analogy holds only in cases where the axis of rotation does not change direction (a fixed axis, rectilinear rolling). These laws may also be applied in non-inertial frames once inertial forces are taken into account.

\subsection{Limits of applicability of Newton's laws of motion}

Newton's laws of motion lose their direct validity in non-inertial frames, where it becomes necessary to introduce inertial forces. Newtonian mechanics also becomes more complicated when coordinate systems other than the traditional Cartesian one are used. % NOTE: the source says "units" (jednostek) where the sentence clearly means "coordinate systems".
In such cases the Lagrangian or Hamiltonian formalisms, based on the principle of least action, work better.

Quantum mechanics describes the motion of particles on the atomic or subatomic scale, where the Newtonian formalism describing the motion of macroscopic bodies does not apply. For example, the uncertainty principle tells us that we cannot simultaneously know the position and the momentum of a particle, whereas classical mechanics rests on their exact trajectories. All of this is, however, a matter of scale: Bohr's correspondence principle tells us that in the limit of large quantum numbers, or of large scales of size or energy, the results of quantum mechanics must agree with the classical ones. In particular, the limiting passage of the Planck constant to $0$ may be regarded as such a transition from a discrete to a continuous system. In that case the Schrödinger equation becomes the classical Hamilton--Jacobi equation.

In his considerations Newton adopted the assumption that time must be absolute, which was only countered 300 years later by Einstein and his special theory of relativity. In that case, when the velocity of the observer is taken into account, the Galilean transformation must be replaced by the Lorentz transformation, in which the Lorentz factor $\gamma=\frac{1}{\sqrt{1 - v^2/c^2}}$ plays an essential role. The resulting phenomena, including those following from general relativity, are again impossible to describe using classical mechanics. They are, however, negligible when $v \ll c$. For this reason classical mechanics still remains a useful tool for describing non-relativistic phenomena.

\section{The postulates of quantum mechanics; measurement, description of the state of a system}

\subsection{The photon hypothesis and its consequences}

\begin{itemize}
    \item \textbf{Light and matter transfer energy to one another in finite portions}
    \begin{equation}
    E = \hbar\omega
    \end{equation}
    \item \textbf{The superposition principle.} If it is possible to prepare a physical system (e.g.\ a particle) both in state A and in state B, then it is also possible to prepare the system under consideration in the superposed state ``A + B'', realising both alternatives at once. The further evolution of the system is linear, and at any later instant the system is in a superposition of the evolved states A and B. In other words: the individual terms of the superposition ``do not know'' about one another.
    \item \textbf{Hilbert space and the representation of the quantum state.} The set of all possible quantum states of any physical system has the structure of a linear vector (Hilbert) space $\mathcal{H}$. To all distinguishable quantum states there correspond basis vectors $\ket{\psi_{k}} \in \mathcal{H}$ of this space. The orthonormality of the basis is given by the scalar product $\braket{\psi_{k}|\psi_{l}} = \delta_{kl}$, where $\bra{\psi_{k}} = \ket{\psi_{k}}^{\dag}$ is the Hermitian conjugate of the basis vector. The vector $\bra{\psi_{k}}$ is referred to as a bra, while the vector $\ket{\psi_{k}}$ is called a ket. The dimension of the Hilbert space $\mathcal{H}$ may be finite or infinite. Superpositions of states correspond to linear combinations of the basis vectors, $\ket{\psi} = \sum_{k}c_{k}\ket{\psi_{k}}$, where $c_{k} \in \mathbb{C}$.
    \item \textbf{Systems with several degrees of freedom.} If a physical system possesses more than one degree of freedom (for instance, a single photon is characterised not only by position but also by polarisation; one may also consider systems composed of a larger number of particles), then to each degree of freedom there corresponds a separate Hilbert space $\mathcal{H}_{i}$, while the Hilbert space of the total system has the structure of a tensor product of the Hilbert subspaces corresponding to the subsystems: $\mathcal{H} =\otimes_{i}\mathcal{H}_{i}$. The dimension of the total Hilbert space is therefore the product, and not the sum, of the dimensions of all the subspaces.
    \item \textbf{Results of physical measurements on quantum states.} Every degree of freedom can be measured. For a system in the state $\ket{\psi} = \sum_{k}c_{k}\ket{\psi_{k}}$ we may, for example, perform a measurement by which we check which of the possibilities $\ket{\psi_{k}}$ is realised. In this situation it is not possible to predict the outcome of the measurement before it is carried out. It is, however, possible to determine the probability of obtaining a given outcome. In the case under consideration, the probability $P_{k}$ of finding the system in the state $\ket{\psi_{k}}$ is $P_{k} = |c_{k}|^{2}$. The complex quantity $c_{k}$ is called the probability amplitude.
\end{itemize}

\subsection{The law of composition of probability amplitudes}

The probability amplitude of a process (for example, one consisting in a photon travelling from the source to the screen) composed of a number of subprocesses (for example, consisting in the photon traversing the path between successive slits) is the product of the probability amplitudes of the individual subprocesses. If, on the other hand, a system starting from some initial state can reach a final state in several ways which are not verified by measurement, then the total amplitude for the system to reach the final state is the sum of the probability amplitudes corresponding to the individual ways of realisation. In other words, if a system may be in one of several possible states and we have no way of determining which is currently realised, then the state of the system is a superposition of all the admissible states.

\subsection{Measurement of a quantum system}

\begin{itemize}
    \item \textbf{Observables.} A physical measurement may give many different outcomes, corresponding to finding the system in one of the possible states $\phi_{k}$. If to each of the possible outcomes we assign some real measured value $a_{k} \in \mathbb{R}$, obtained with probability $P_{k} = |\braket{\phi_{k}|\psi}|^{2}$, then the mean measurement outcome is obtained from the classical formula of probability theory: % NOTE: the source has a malformed bra-ket here; corrected to |<phi_k|psi>|^2.
    (mean measurement outcome) $= \sum_{k}a_{k}P_{k} = \sum_{k}a_{k}|\braket{\phi_{k}|\psi}|^{2} = \bra{\psi}\left(\sum_{k}a_{k}\ket{\phi_{k}}\bra{\phi_{k}}\right)\ket{\psi}$, where $\ket{\psi_{k}}\bra{\psi_{k}}$ is the projection operator onto the state $\ket{\psi_{k}}$. To every possible observable, that is, to every physical quantity that can be measured on the system and that leads exclusively to real measurement outcomes, there corresponds some Hermitian operator $\hat{A} = \hat{A}^{\dag}$. In the case of an observable giving the possible numerical outcomes $a_{k}$, corresponding to finding the system in the corresponding states $\ket{\psi_{k}}$, this operator is:
    \begin{equation}
        \hat{A} = \sum_{k}a_{k}\ket{\phi_{k}}\bra{\phi_{k}}
    \end{equation}
    The mean value of the observable so defined takes the form:
    \begin{equation}
        \langle\hat{A}\rangle = \braket{\psi|\hat{A}|\psi}
    \end{equation}
\end{itemize}

\subsection{Time evolution of the quantum state}

\begin{itemize}
    \item \textbf{Unitarity of the evolution.} Consider two distinguishable, and hence orthogonal, quantum states $\psi$ and $\phi$ prepared at the initial instant. We adopt as a postulate the thesis that under the passage of time these states will remain fully distinguishable. Since the degree of distinguishability of two states is given by the amplitude $\braket{\phi|\psi}$, we shall assume in general that the evolution of quantum states does not change the scalar products between the corresponding vectors.
    \item \textbf{Evolution equation of the quantum state vector}
    \begin{equation}
        i\hbar\frac{d}{dt}\ket{\psi\left(t\right)} = \hat{H}\ket{\psi\left(t\right)}
    \end{equation}
    In the case of a time-independent Hamiltonian:
    \begin{equation}
        \hat{H}\ket{\psi\left(r\right)} = \mathsf{E}\ket{\psi\left(r\right)}
    \end{equation}
    \begin{equation}
        \ket{\psi\left(r,t\right)} = e^{-\frac{i}{\hbar}\mathsf{E}t}\ket{\psi\left(r\right)}
    \end{equation}
\end{itemize}

\subsection{The wave function}

The wave function is a mathematical description of the quantum state. It takes complex values, and its squared modulus has the interpretation of a probability density over the space on which it is defined.

% NOTE: the source has \int_{b}^{b}; the lower limit should be a.
\begin{equation}
    P_{a \le x \le b}\left(t\right) = \int^{b}_{a}|\Psi(x,t)|^{2}dx
\end{equation}

The wave function satisfies a normalisation condition:

% NOTE: the source has an extra "dt" in the integrand; the normalisation is over x only.
\begin{equation}
    \int^{\infty}_{-\infty}|\Psi(x,t)|^{2}dx = 1
\end{equation}

Knowing the form of the wave function, the quantum state vector may be determined:

\begin{equation}
    \ket{\Psi\left(t\right)} = \int \Psi\left(x,t\right)\ket{x}dx
\end{equation}

\section{Angular momentum and the phenomenon of precession --- classical and quantum description}

The classical angular momentum is defined as:

\begin{equation}
    \vec{L} = \sum_i \vec{r}_i \times \vec{p}_i
\end{equation}

where $r$ is the position and $p$ the momentum of the $i$-th element of the system.

In quantum mechanics the orbital angular momentum is defined as an operator composed of three components and given by the formula

\begin{equation}
    \hat{L} = -i\hbar(\hat{r}\times \nabla)
\end{equation}

A further operator corresponding to angular momentum is the spin angular momentum operator $\hat{S}$, which likewise consists of three components, being the projections onto the respective axes. The total angular momentum $\hat{J}$ is the sum of the spin and orbital angular momenta, i.e.\ $\hat{J} = \hat{L} + \hat{S}$. Quantisation permits the eigenvalues of the squares of the angular momentum operators and of their projections onto the $z$ axis to be found, i.e.\ $\hat{L}^2\ket{l} = \hbar^2 l(l+1)\ket{l}$, $\hat{L_z}\ket{m} = \hbar m\ket{m}$. Quantum numbers for $\hat{S}$ and $\hat{J}$ may be defined analogously.

Precession of a rigid body occurs when the axis about which the body rotates does not coincide with any of the principal axes of the inertia tensor. Forced precession occurs when a torque perpendicular to the angular momentum acts on a rotating body. The axis of rotation then executes a motion tracing out a surface in the shape of the lateral surface of a cone. This phenomenon may be observed in the example of a spinning top.

Classically, the precession frequency is given by the formula

\begin{equation}
    \omega_p = \frac{\tau}{I_s\omega_s\sin(\theta)}
\end{equation}

% NOTE: the source does not define omega_s; it is the spin angular velocity of the body.
where $I_s$ is the moment of inertia of the object, $\omega_s$ its spin angular velocity, $\tau$ the torque and $\theta$ the angle between the axis of rotation and the axis of precession.

In quantum mechanics one of the corrections to the spin--orbit interaction is the phenomenon of Thomas precession, which takes into account relativistic time dilation. It describes the phenomenon whereby objects undergo precession when they accelerate in special relativity. Larmor precession is a phenomenon exploited in studying the interaction of electrons with a magnetic field. It is the precession of the magnetic moment of an electron in an external magnetic field. This phenomenon is used in magnetic resonance, because the precession frequency does not depend on the spatial orientation of the spins.

The torque induced by a magnetic field $\vec{B}$ is given by the formula

\begin{equation}
    \vec{\mu}\times \vec{B}
\end{equation}

where $\vec{\mu}$ is the magnetic dipole moment.

\section{Variational principles of physics; examples of equations that can be derived from them}

Variational principles allow the equations of motion to be formulated by minimising a certain function called the action functional. The basic idea is that the actual trajectory of a physical system is the one that minimises (or renders stationary) the action $S$. A functional is a certain mapping from a vector space into $\mathbb{R}$, and the task of the calculus of variations is to find the extremum of such an action. The considerations born of variational principles created a tool which turned out to be easy to adapt to the description of quantum phenomena, and it is precisely in them that the theoretical foundations for their application may be sought. The trajectories that may be taken are associated with different probability amplitudes.

The aforementioned idea is called the principle of least action and has its beginnings as far back as antiquity, when it was stated that light travels along a trajectory that minimises distance and time simultaneously. This principle was carried over into the field of classical mechanics, in part from the brachistochrone problem. The puzzle formulated by Bernoulli consisted in determining the curve along which a rolling bead would reach a lower-lying point most quickly and by the shortest route. Years later, the particular solution led Lagrange to create a general analytical formalism permitting the equations of motion to be written in the calculus of variations, and the trajectories then to be determined. The action is here defined as the integral over time of the Lagrangian $\mathcal{L}$, which is the difference of the kinetic energy $T$ and the potential energy $V$:

\[
S = \int_{t_1}^{t_2}\mathcal{L}dt = \int_{t_1}^{t_2}(T-V)dt
\]

The principle of least action assumes that the action is stationary, that is, just as in the case of the extremum of a function, infinitesimally small deviations will not change the result. From this condition we obtain the variational principle for classical mechanics in the form of the Euler--Lagrange equation:

\[
\frac{d}{dt}\left(\frac{\partial \mathcal{L}}{\partial \dot q_i}\right) - \frac{\partial\mathcal{L}}{\partial q_i} = 0
\]

where $q_i$ are the generalised coordinates and $\dot q_i$ the corresponding generalised velocities. The formalism created by Lagrange works excellently for calculating trajectories in complicated systems, especially on account of the generalised variables. By means of a Legendre transformation, on the other hand, one may pass over to Hamiltonian mechanics. While in the context of classical mechanics both formalisms describe essentially the same phenomena, the latter has far more applications. The Hamiltonian, as a measure of the total energy of the system, also permits a convenient description of problems in statistical physics, and even in quantum mechanics. Applying the calculus of variations and representing the motion of a particle as a wave, one may formulate the Hamilton--Jacobi equation, which is very close in idea to the Schrödinger equation:

\[
H \left( \vec{q}, \frac{\partial S}{\partial \vec{q}}, t \right) + \frac{\partial S}{\partial t} = 0
\]

% NOTE: the source writes the Lagrangian as -\frac{1}{4}F_{\mu\nu}^{\mu\nu};
% the correct form is -\frac{1}{4}F_{\mu\nu}F^{\mu\nu}.
The applications of the variational principle are not, however, limited to classical mechanics. For example, it is also possible to obtain Maxwell's equations in tensor form by applying it to the Lagrangian of the electromagnetic field in vacuum, $\mathcal{L} = -\frac{1}{4}F_{\mu\nu}F^{\mu\nu}$. Moreover, it is possible to obtain the full Einstein equations by applying the principle to the Hilbert action.

\section{The laws of thermodynamics}

\subsection{The zeroth law of thermodynamics}

If systems A and B, capable of exchanging heat with one another, are in thermal equilibrium with each other, and the same is true of systems B and C, then systems A and C are also in thermal equilibrium with each other.

\subsection{The first law of thermodynamics}

The change in the internal energy of a system is equal to the sum of the heat supplied and the work done on the system.

\begin{equation}
    \Delta\mathsf{U} = \mathsf{W} + \mathsf{Q}
\end{equation}

\subsection{The second law of thermodynamics}

\begin{itemize}
    \item Clausius statement: \textit{There is no thermodynamic process whose sole result would be the absorption of heat from a reservoir at a lower temperature and its transfer to a reservoir at a higher temperature.}
    \item Kelvin statement: \textit{No process is possible whose sole effect would be the absorption of a certain quantity of heat from a reservoir and its conversion into an equivalent amount of work.}
    \item Ostwald statement: \textit{There is no perpetual motion machine of the second kind (a heat engine taking thermal energy from a system and converting it entirely into work).}
\end{itemize}

\subsection{The third law of thermodynamics}

The Nernst heat theorem: \textit{It is not possible, by means of a finite number of steps, to attain the temperature of absolute zero (zero kelvin) if we take a non-zero absolute temperature as our starting point.}

\section{Mechanical properties of material media: liquids, elastic media, mechanical waves}

Liquids are substances intermediate between a gas and a solid, having an easily variable shape but little compressibility. The properties of a liquid follow from the behaviour of its molecules:

\begin{itemize}
    \item they have complete freedom of movement within the volume occupied by the liquid
    \item intermolecular interactions occur between them, which within the bulk of the liquid cancel one another out
    \item the intermolecular interactions do not cancel at the interface between the liquid and another phase, as a result of which the phenomenon of surface tension occurs.
\end{itemize}

An elastic medium is a medium in which an external force gives rise to stresses whose effect is the appearance of an outwardly directed force. This force causes a body forming part of that medium to return to its original shape. Strain is understood as a change in the size of the medium.

While most solids are elastic over a very wide range of stresses, the elasticity of fluids (liquids and gases) depends on several conditions. If a fluid is enclosed in an elastic container, e.g.\ air in a balloon, such a system is elastic. If an external force acts at high frequency on a fluid in the free state, then such a system is also elastic. In this way sound propagates in air as a mechanical wave.

\paragraph{Mechanical waves}

If we displace some fragment of an elastic medium from its equilibrium position, it will consequently oscillate about that position. These oscillations, thanks to the elastic properties of the medium, are transmitted to successive parts of the medium, which begin to oscillate. In this way the disturbance passes through the whole medium. The speed of propagation of waves depends on the properties of the medium. According to the shape of the wavefront we may distinguish plane waves and spherical waves.

According to the direction of oscillation of the particles of the medium relative to the propagation of the waves, we may distinguish transverse waves and longitudinal waves.

A wave is longitudinal when the direction of oscillation of the particles of the medium is parallel to the direction of propagation of the wave and at the same time to the direction of energy transport. Examples here are sound waves in air, or the oscillations of a spring alternately compressed and extended.

A wave is transverse when the direction of oscillation of the particles of the medium is perpendicular to the direction of propagation of the wave and at the same time to the direction of energy transport. An example here may be the oscillations of a taut string one end of which we move cyclically up and down.

\section{Flow of an incompressible liquid; Bernoulli's principle}

Let us begin by defining the assumptions of the problem. It is an example from the field of ideal-fluid hydrodynamics, and therefore adopts certain simplifications:

\begin{itemize}
    \item the fluid is incompressible, that is, it does not change its volume under pressure, which for liquids is usually a very good approximation
    \item friction arising from the viscosity of the liquid is negligible
    \item the flow of the liquid is stationary and irrotational, which amounts to its parameters (velocity, density, \ldots) being independent of time
\end{itemize}

Applying the principle of conservation of energy to the flow problem, one obtains Bernoulli's principle, which is usually written as:

\[
\frac{v^2}{2} + gz + \frac{p}{\rho} = const
\]

% NOTE: the source's explanation of the hydrodynamic paradox is garbled and
% repeats the term twice; retained as written.
where we have, in turn, the kinetic energy, the potential energy and the energy due to pressure of a unit of fluid. For a chosen piece of liquid this equation will always be satisfied, and it permits the description of flow through a pipe changing its cross-section or height. For example, one may demonstrate the so-called hydrodynamic paradox by analysing a widening pipe. For a smaller cross-section the liquid flows faster, which must be compensated by a lower pressure over the shorter section. This counter-intuitive property is also known as the hydrodynamic paradox.\merr

\section{Maxwell's equations, the scalar and vector potentials of the electromagnetic field, the equations of the electromagnetic field in vacuum and in material media}

\subsection{Maxwell's equations in vacuum}

\begin{equation}
    \begin{matrix}
    \nabla \vec{E} = \frac{\rho}{\epsilon_{0}}\\
    \nabla \vec{B} = 0\\
    \nabla\times\vec{E} = -\frac{\partial \vec{B}}{\partial t}\\
    \nabla\times\vec{B} = \mu_{0}\vec{J} + \mu_{0}\epsilon_{0}\frac{\partial\vec{E}}{\partial t}
    \end{matrix}
\end{equation}

\subsection{Maxwell's equations in terms of free currents and charges}

We define the electric displacement field $\vec{D}$ and the magnetic field strength $\vec{H}$ as:

\begin{equation}
    \begin{matrix}
        \vec{D} = \epsilon_{0}\vec{E} + \vec{P}\\
        \vec{H} = \frac{1}{\mu_{0}}\vec{B} - \vec{M}
    \end{matrix}
\end{equation}

Maxwell's equations may then be rewritten as

\begin{equation}
    \begin{matrix}
    \nabla \vec{D} = \rho_{f}\\
    \nabla \vec{B} = 0\\
    \nabla\times\vec{E} = -\frac{\partial \vec{B}}{\partial t}\\
    \nabla\times\vec{H} = \vec{J}_{f} + \frac{\partial\vec{D}}{\partial t}
    \end{matrix}
\end{equation}

\subsection{The vector potential of the electromagnetic field}

The magnetic vector potential $\vec{A}$ is a three-dimensional vector field whose curl is the magnetic field.

\begin{equation}
    \vec{B} = \nabla\times\vec{A}
\end{equation}

\subsection{The scalar potential of the electromagnetic field}

For time-dependent potentials the electric field may be written in the form:

\begin{equation}
    \vec{E} = -\nabla\Phi - \frac{\partial\vec{A}}{\partial t}
\end{equation}

where $\Phi$ is the electric potential.

\section{The electric field in material media; dielectric polarisation, conductivity, mobility of charge carriers}

\subsection{Polarisation}

Dielectric polarisation is the phenomenon consisting in the formation of electric dipoles, or in the orientation of existing ones, under the influence of an external electric field. As a result of polarisation an internal electric field arises which partially balances the external field. The polarisation vector arising as a result of the action of the external field may, in the linear approximation, be written as:

\begin{equation}
    \vec{P} = \varepsilon_0\chi\vec{E}
\end{equation}

Dielectric polarisation is the result of several basic phenomena:

\paragraph{Electronic polarisation}

Electronic polarisation is the result of the displacement, under the influence of an electric field, of the negatively charged electron clouds with respect to the positive nuclei. It occurs in all materials. It is a very fast process, and therefore, when investigated with an alternating electric field, the effects associated with it are observed up to frequencies corresponding to the ultraviolet. The dipole moment $p_e$ is determined by the electronic polarisability $\alpha_e$:

\begin{equation}
    p_e = \alpha_e \cdot E
\end{equation}

\paragraph{Ionic polarisation}

Ionic polarisation occurs only in materials with ionic bonds. It is the result of the displacement of ions in the crystal lattice of the material---those of one sign in one direction, those of the other sign in the other. It is the slowest of all the polarisation processes, since it requires the movement of many atoms bound in the crystal lattice.

\paragraph{Orientational polarisation}

Orientational polarisation, also called dipolar polarisation, occurs only in polar dielectrics, that is, those whose molecules form permanent dipoles (they possess their own dipole moment). In an electric field an ordering torque acts on them. At the same time the ordering is destroyed by thermal vibrations, which causes orientational polarisation to be temperature dependent.

\subsection{Conductivity and mobility}

Conduction is the phenomenon of the flow of electric charges occurring in a material under the influence of an electric field. The quantity characterising electrical conduction is the conductivity $\sigma$:

\begin{equation}
    \sigma = q n \mu
\end{equation}

where $q$ is the charge of the carriers, $n$ their concentration and $\mu$ the carrier mobility. The mobility of the carriers depends on the applied field and on the drift velocity. Temperature is one of the factors affecting conductivity. A change in temperature may affect a change both in the concentration and in the mobility of the carriers. The most common types of conduction are: electronic, where the charge carriers are electrons, hole conduction, where the charge carriers are electrons moving in the valence band, and ionic, where the charge carriers are ions.

\section{The magnetic field in material media; classification of magnetic media, microscopic explanations of the exchange interaction between the magnetic moments of the constituents (atoms, ions, molecules)}

On the microscopic scale every medium exhibits magnetic properties, on account of the spin of its constituent particles and of the motion of electrons about the atomic nuclei. On the macroscopic scale, however, magnetic phenomena are by no means so common, which follows from the absence of long-range order and from destructive interference between the constituent magnetic moments. Long-range order may nevertheless be introduced as a result of an external magnetic field $B$, similarly to the case of inducing polarisation of a medium with an electric field. In general, when describing magnetic fields in material media we use constitutive relations analogous to those for the electric field:

\begin{align*}
    \vec{H} &= \frac{1}{\mu_0}\vec{B} - \vec{M} \\
    \vec{M} &= \chi_m\vec{H} \\
    \vec{H} &= \frac{1}{\mu_0(1+\chi_m)}\vec{B}
\end{align*}

where $\vec{H}$ is the magnetic field strength, $\vec{M}$ is the magnetisation of the medium and $\chi_m$ the magnetic susceptibility. Interactions with a magnetic field are, however, more nuanced than in the case of the electric field. According to the Bohr--van Leeuwen theorem, a classical system will always have zero mean magnetisation, and therefore any magnetisation must be the result of quantum effects. Three kinds of medium are usually distinguished.

\subsection{Paramagnets}

The most common group, characterised by positive magnetic susceptibility and hence by a magnetisation whose direction coincides with that of the magnetic field strength. It arises from the presence of unpaired electrons which, under the influence of an external field, align their magnetic moments with its direction. The most commonly considered models of paramagnetism are:

\begin{itemize}
    \item Curie paramagnetism --- considering a system of isolated magnetic moments, one may derive a formula for the susceptibility of the whole medium as $\chi_m = \frac{C}{T}$.
    \item Pauli paramagnetism --- carriers moving freely inside a metal may be treated as a Fermi gas, and their interactions with an external field lead to macroscopic magnetisation.
\end{itemize}

\subsection{Diamagnets}

Diamagnetism is the phenomenon of inducing a magnetisation opposite to the magnetic field strength, which may be expressed by a negative susceptibility. It occurs in every material, but is usually weaker than paramagnetism, and therefore we classify as diamagnets those media in which it is the dominant effect. Thus the condition is usually that all electrons be paired. A special case is that of superconductors, which expel the magnetic field completely through the Meissner effect, while the remainder are described by:

\begin{itemize}
    \item Langevin diamagnetism --- describes the diamagnetism of dielectrics, where electrons under the influence of an external field begin to circulate, which is called Larmor precession. A charge moving in a circle is de facto a loop through which a current flows, which effectively gives a magnetic moment of opposite direction.
    \item Landau diamagnetism --- similarly, it considers the Fermi gas in metals, except that in this case the orbital moment, and not the spin of the particles, is taken into account.
\end{itemize}

\subsection{Ferromagnets}

Taking into account the exchange interactions between microscopic magnetic moments leads to a third category, ferromagnetism. Here the particles tend to align not only with the external field but also with their neighbours, leading to the formation of magnetic domains within the material, resulting in a strong magnetic susceptibility and also in the occurrence of remanent magnetisation. The analysis of such materials is carried out by means of statistical physics, and depending on the material one may make use of the Ising, XY or Heisenberg models. The most general is the last of these, which permits arbitrary alignments of the moments, while the Ising model distinguishes only up and down, and the XY model rotation about a single axis.

It is also worth mentioning the existence of antiferromagnets and ferrimagnets. In contrast to ferromagnets, their domains are not aligned in one direction. For the former they cancel completely, and for the latter only partially.

\section{Electromagnetic waves as examples of solutions of Maxwell's equations; polarisation, reflection and refraction of waves, the Doppler effect}

\subsection{Electromagnetic waves as examples of solutions of Maxwell's equations}

Let us recall the general Maxwell equations:

\begin{equation}
    \begin{matrix}
    \nabla \vec{E} = \frac{\rho}{\epsilon_{0}}\\
    \nabla \vec{B} = 0\\
    \nabla\times\vec{E} = -\frac{\partial \vec{B}}{\partial t}\\
    \nabla\times\vec{B} = \mu_{0}\vec{J} + \mu_{0}\epsilon_{0}\frac{\partial\vec{E}}{\partial t}
    \end{matrix}
\end{equation}

We shall find an electromagnetic wave that is a solution of this system of equations in vacuum, that is, when $\rho = 0$ and $\vec{J} = 0$. The equations then take the following form:

\begin{equation}
    \begin{matrix}
    \nabla \vec{E} = 0\\
    \nabla \vec{B} = 0\\
    \nabla\times\vec{E} = -\frac{\partial \vec{B}}{\partial t}\\
    \nabla\times\vec{B} = \mu_{0}\epsilon_{0}\frac{\partial\vec{E}}{\partial t}
    \end{matrix}
\end{equation}

We take the curl of the third and fourth Maxwell equations. Consider the third equation:

\begin{equation}
    \nabla\times\left(\nabla\times\vec{E}\right) = -\nabla\times\left(\frac{\partial\vec{B}}{\partial t}\right)
\end{equation}

We make use of the $\vec{BAC}-\vec{CAB}$ identity:

\begin{equation}
    \vec{A}\times\left(\vec{B}\times\vec{C}\right) = \vec{B}\cdot\left(\vec{A}\cdot\vec{C}\right) - \vec{C}\cdot\left(\vec{A}\cdot\vec{B}\right)
\end{equation}

The left-hand side of the third Maxwell equation:

\begin{equation}
    \mathrm{LHS} = \nabla\times\left(\nabla\times\vec{E}\right) = \nabla\left(\nabla\cdot\vec{E}\right) - \nabla^{2}\vec{E} = -\Delta\vec{E}
\end{equation}

The right-hand side of the third Maxwell equation:

\begin{equation}
    \mathrm{RHS} = -\nabla\times\left(\frac{\partial\vec{B}}{\partial t}\right) = - \frac{\partial}{\partial t}\left(\nabla\times\vec{B}\right) = -\mu_{0}\epsilon_{0}\frac{\partial^{2}\vec{E}}{\partial t^{2}}
\end{equation}

Equating both sides we obtain:

\begin{equation}
    \left(\Delta - \mu_{0}\epsilon_{0}\frac{\partial^{2}}{\partial t^{2}}\right)\vec{E} = 0
\end{equation}

We proceed analogously with the fourth Maxwell equation and find the equation for $\vec{B}$.

\begin{equation}
    \left(\Delta - \mu_{0}\epsilon_{0}\frac{\partial^{2}}{\partial t^{2}}\right)\vec{B} = 0
\end{equation}

From the above two equations it follows that the speed of propagation of an electromagnetic wave in vacuum is $c = \frac{1}{\sqrt{\mu_{0}\epsilon_{0}}}$.

\subsection{Polarisation of a wave}

Polarisation is a property of a transverse wave concerning the ordered relation between the direction of oscillation of the disturbance and the direction of propagation of the wave. The best known polarisations are linear and circular. Polarisation occurs only for those kinds of wave, and under those conditions, in which the oscillations may take place in various directions perpendicular to the direction of propagation of the wave.

\subsection{The laws of reflection and refraction of electromagnetic waves}

\textbf{The law of reflection.} The angle of reflection is equal to the angle of incidence, and the incident ray, the reflected ray and the normal to the reflecting surface lie in one plane. As a result of reflection only the direction of propagation of the wave changes, and its wavelength does not change.

\textbf{The law of refraction.} The law of refraction (Snell's law) states that the incident and refracted rays and the normal to the surface lie in one plane, and that the angles satisfy the relation:

\begin{equation}
    \frac{\sin{\theta_{1}}}{\sin{\theta_{2}}} = \frac{n_{2}}{n_{1}}
\end{equation}

where $n_{1}$ is the refractive index of the first medium, $n_{2}$ the refractive index of the second medium, $\theta_{1}$ the angle of incidence, that is, the angle between the incident ray and the normal to the interface between the media, and $\theta_{2}$ the angle of refraction, the angle between the refracted ray and the normal.

\textbf{The Brewster angle} --- the angle of incidence of light on the surface of a dielectric at which the reflected ray is completely linearly polarised. When unpolarised light falls on the boundary between transparent media at such an angle that the reflected and refracted rays form a right angle, the reflected light is completely polarised. The Brewster angle is given by the formula:

\begin{equation}
    \tan{\alpha_{B}} = \frac{n_{2}}{n_{1}}
\end{equation}

where $n_{1} < n_{2}$, and the light is incident on the interface from the side of medium $n_{1}$.

\subsection{The Doppler effect}

\textbf{The Doppler effect} --- a physical phenomenon occurring for waves, consisting in the appearance of a difference between the frequency of the wave emitted by its source and the frequency of the wave recorded by an observer who is moving with respect to that source. For waves propagating in a medium, such as for example sound waves, the magnitude of the effect depends on the velocities of the observer and of the source relative to the medium in which these waves propagate. In the case of waves propagating without the participation of a material medium, such as for example light (electromagnetic waves in general) in vacuum, only the difference of the velocities of source and observer matters.

\textbf{Acoustic waves}

In the case where the observer is moving, the wave speed is effectively changed

\begin{equation}
    f' = \frac{V_{o} + V}{\lambda} = \frac{V}{\lambda}\frac{V_{o} + V}{V} = f\frac{V_{o} + V}{V}
\end{equation}

where $f$ is the frequency of the wave emitted by the source and $V_{o}, V$ are the velocities of the observer and of the wave. In this convention an observer approaching the source has a velocity of positive sign. In the case where the source is moving, the wavelength emitted is effectively changed:

\begin{equation}
    f' = \frac{V}{\lambda'} = \frac{V}{\lambda - V_{z}T} = \frac{V}{VT - V_{z}T} = f\frac{V}{V - V_{z}},
\end{equation}

where $T$ is the period of the wave and $V_{z}$ is the velocity of the source. In this convention a source approaching the observer has a velocity of positive value.

\textbf{Electromagnetic waves}

Here the effect is already a relativistic one. We take into account only the relative velocity, i.e.\ it does not matter whether the source (or the observer) is in motion. On passing from one frame of reference to another, for a source or observer receding from one another, the frequency of the wave is:

\begin{equation}
    f' = f\frac{\sqrt{1-\beta^{2}}}{1+\beta}
\end{equation}

where $\beta = \frac{v}{c}$. If the motion takes place at an angle $\phi$, then:

\begin{equation}
    f' = f\frac{\sqrt{1-\beta^{2}}}{1+\beta\cos{\phi}}
\end{equation}

For the angle $\phi = \frac{\pi}{2}$ the so-called transverse Doppler effect occurs, which does not arise for mechanical waves.

\section{Mechanical waves, electromagnetic waves, matter waves; similarities and differences}

Electromagnetic waves are described by the equation:

\begin{equation}
    \nabla^2\vec{E} = \mu_0 \varepsilon_0\frac{\partial^2}{\partial t^2}\vec{E}
\end{equation}

For a plane wave propagating in the $z$ direction with angular frequency $\omega$ and wave number $k$, one of the solutions of the above equations is:

\begin{equation}
    E(z,t) = E_0\cos \left( \omega t - kz  \right)
\end{equation}

In quantum mechanics, electromagnetic radiation is described as a stream of photons---wave packets. The energy of a single photon depends on the oscillation frequency of the electromagnetic wave, $E = \hbar \omega$, while the momentum is described as $p = \hbar k$.

Matter waves --- every material object may be described in two ways: as a collection of particles or as a wave. Effects confirming the wave nature of matter are observed in the form of the diffraction of elementary particles, and even of whole atomic nuclei. The formula for the de Broglie wavelength is $\lambda = \frac{h}{mv}$. It is relatively easy to observe wave effects in the case of light particles, e.g.\ electrons (small objects manifest wave properties). The diffraction and interference of electron waves are exploited in electron diffraction imaging and in low-energy electron diffraction. The above considerations concern the free motion of particles (to which plane waves would correspond). In real cases a certain group of matter waves, a so-called wave packet, must be assigned to the particle.

Mechanical waves --- if we displace some fragment of an elastic medium from its equilibrium position, it will consequently oscillate about that position. These oscillations, thanks to the elastic properties of the medium, are transmitted to successive parts of the medium, which begin to oscillate. In this way the disturbance passes through the whole medium. The speed of propagation of waves depends on the properties of the medium.

All waves are subject to typically wave-like phenomena such as interference and diffraction, and satisfy the laws of reflection and refraction. As a result of superposition a standing wave may arise.

Electromagnetic waves do not need a medium in order to propagate, i.e.\ they can propagate in vacuum.

\section{Phase transitions; kinds of phase transition, the order parameter of the transition, conditions for the coexistence of phases (examples of phenomena)}

% NOTE: two points below follow the source but are questionable:
% (i) the Landau description uses a free energy, not the enthalpy;
% (ii) the order parameter conventionally takes larger (non-zero) values in
% the lower-entropy, ordered phase --- the source states the opposite.
A substance which is physically and chemically homogeneous throughout its volume is called a phase. With a change in external factors, such as temperature or pressure, an abrupt change in these properties may however occur, which is termed a phase transition. These changes may concern the state of matter, the crystal structure, magnetic properties or superconductivity. In the phenomenological description of thermodynamics, phase transitions are described by means of an enthalpy\merr{} dependent on the order parameter of the transition. Depending on the case, this is a dynamical variable which takes lower values for the phase of lower entropy\merr, i.e.\ for example the magnetisation for ferromagnets or the density for liquids.

The aforementioned discontinuous changes of the parameters of the system at a phase transition lead to the Ehrenfest classification, in which changes of the chemical potential $\mu$ are considered. By a phase transition of the $n$-th kind we understand one in which the lowest partial derivative $\frac{\partial^n \mu}{\partial T^n}$ is discontinuous. Examples of first-order phase transitions are changes of the state of matter, where a discontinuous change of density occurs. An example of the second order, on the other hand, is that of ferromagnets, which below the Curie temperature begin to increase their magnetic susceptibility. In practice, third-order phase transitions do not occur. This classification also allows for continuous phase transitions, an example of which may be the superfluid state.

The conditions for the coexistence of phases are given by the Gibbs phase rule. It states what the number of intensive parameters $f$ is that may be varied without changing the number of phases in equilibrium, and is given by:

\[ f = s - r + 2 \]

where $s$ is the number of chemical substances and $r$ the number of phases. For the rule to be satisfied, the system must be in thermal, mechanical and chemical equilibrium. For example, for a single substance, $s=1$, we have:

\begin{itemize}
    \item $r=1$, one phase, i.e.\ $f=2$ --- a surface
    \item $r=2$, two phases at once, $f=1$ --- a curve
    \item $r=3$, three phases, $f=0$ --- a (triple) point
\end{itemize}

\section{Experimental grounds for the emergence of quantum mechanics; examples of phenomena incomprehensible on classical grounds and their explanations within quantum theory}

\subsection{The photoelectric effect}

The photoelectric effect is the emission of electrons from a material caused by electromagnetic radiation (light). Electrons emitted in this way are called photoelectrons. The quantum hypothesis explains why the energy of the photoelectrons depends on the frequency of the light, and why below a certain frequency of light the photoelectric effect does not occur.

\subsection{The Compton effect}

The phenomenon of the scattering of X-radiation and gamma radiation, that is, of electromagnetic radiation of high frequency, on free or weakly bound electrons, as a result of which the wavelength of the radiation increases. The scattering is caused by collisions of individual quanta of radiation with electrons.

\subsection{Young's experiment}

An experiment consisting in passing coherent light through two closely spaced slits and observing the pattern formed on a screen. As a result of interference, bright and dark fringes appear on the screen in the regions where the light is cancelled or reinforced. Other objects on the atomic scale, such as electrons, exhibit the same behaviour when they are fired towards a double slit. Moreover, it has been observed that the detection of the individual discrete impacts is inherently probabilistic, which is inexplicable using classical mechanics.

\subsection{Quantisation of the energy levels in the hydrogen atom}

Atomic spectroscopy shows that there exists a discrete infinite set of states in which the hydrogen atom (or any other) may exist, contrary to the predictions of classical physics. Attempts at a theoretical understanding of the states of the hydrogen atom were important for the history of quantum mechanics, because all other atoms can be roughly understood once this simplest atomic structure is known in detail.

\subsection{The ultraviolet catastrophe}

The ultraviolet catastrophe, also called the Rayleigh--Jeans catastrophe, was a prediction of classical physics from the end of the nineteenth and the beginning of the twentieth century, that an ideal black body in thermal equilibrium would emit an unbounded amount of energy as the wavelength fell into the ultraviolet range. In 1900 Max Planck derived the correct form of the spectral intensity distribution function by adopting certain assumptions that were strange (by the standards of the time). In particular, Planck assumed that electromagnetic radiation may be emitted or absorbed only in discrete packets of energy, called quanta.

\section{The uncertainty principle and the limits of measurement accuracy}

% NOTE: as written in the source, the left-hand side is a product of variances
% while the right-hand side is not squared; the standard Robertson relation is
% $\Delta^2\hat A\,\Delta^2\hat B \geq \frac{1}{4}|\langle[\hat A,\hat B]\rangle|^2$.
The Heisenberg uncertainty principle states that there exist pairs of quantities which cannot be measured simultaneously to arbitrary accuracy. Let $\hat{A}$, $\hat{B}$ be observables and $\Delta^2 \hat{A} = \bra{\psi}\hat{A}^2\ket{\psi} - \left|\bra{\psi}\hat{A}\ket{\psi}\right|^2$. Then:

\begin{equation}
    \Delta^2 \hat{A} \Delta^2 \hat{B} \geq \frac{1}{4}\left| \langle [\hat{A},\hat{B}] \rangle  \right| \qquad \merr
\end{equation}

In particular, $\Delta^2 x \cdot \Delta^2 p_x \geq \frac{\hbar^2}{4}$. The state that saturates this inequality is a Gaussian wave packet. This means that there exist no states guaranteeing better localisation of the particle in position and momentum space.

The uncertainty principle should not be confused with ``Heisenberg's microscope'', which states that in order to measure the position of a particle ever more accurately, a stream of photons of ever shorter wavelength must be used. However, the photons used for illumination consequently have ever greater momentum and in consequence cause a greater uncertainty of momentum.

It is worth noting that such reasoning does not capture the meaning of the uncertainty principle, since it assumes the existence of an objective position of the electron which can be measured.

\section{Approximate methods of quantum mechanics; perturbation theory for eigenvalues, Fermi's golden rule, the Born approximation for scattering amplitudes}

Quantum mechanics describes the behaviour of systems at the atomic and subatomic level. However, it rests to a large extent on solving differential equations which we are able to solve analytically only in relatively simple cases. Sometimes even determining the Hamiltonian itself is difficult, and therefore for more complex systems approximate methods are used, which make it possible to obtain useful results.

\subsection{Stationary perturbation theory}

Given a known Hamiltonian $H_0$, for example that of the hydrogen atom, we may wish to introduce a slight modification into the system, for example a weak external electric field. Solving it analytically again would be just as highly non-trivial as the hydrogen atom itself. We may instead treat our correction as a perturbation in the form of an additional term $\lambda H'$. The method itself is the quantum counterpart of an expansion in a Taylor series with respect to the parameter $\lambda$ and of calculating successive corrections. The smaller the perturbation, the fewer corrections need to be calculated, and often one or two suffice. They are given by the following formulae:

\begin{align*}
    E_n^{(1)} &= \bra{\psi_n^{(0)}} H' \ket{\psi_n^{(0)}} \\
    \psi_n^{(1)} &= \sum_{m \neq n} \frac{\bra{\psi_m^{(0)}} H' \ket{n^{(0)}}}{E_n^{(0)} - E_m^{(0)}} \psi_m^{(0)} \\
    E_n^{(2)} &= \sum_{m \neq n} \frac{|\bra{\psi_m^{(0)}} H' \ket{n^{(0)}}|^2}{E_n^{(0)} - E_m^{(0)}}
\end{align*}

\subsection{Fermi's golden rule}

The method of stationary perturbation theory described above assumes that the perturbing potential is invariant in time. There also exists a dynamical method, but it is far more complicated. In the special case, however, in which the perturbing potential varies periodically in time, one may apply Fermi's golden rule, an important result of that calculation. This solution gives the probability per unit time of transitions between two quantum states $\psi_n^{(0)}$ and $\psi_m^{(0)}$, for a perturbation $H'e^{-i\omega t}$:

\[
\Gamma_{n\rightarrow{}m} \approx \frac{2\pi}{\hbar} |\bra{\psi_m^{(0)}}H'\ket{\psi_n^{(0)}}|^2 \delta\left(E_n^{(0)} - E_m^{(0)} - \hbar\omega\right)
\]

\subsection{The Born approximation}

% NOTE: the source's prefactor is 2m/(K\hbar); the standard first Born result
% for a spherically symmetric potential carries \hbar^2 and a minus sign:
% f(\theta) = -\frac{2m}{K\hbar^2}\int_0^\infty r'V(r')\sin(Kr')dr'.
A further important result of perturbation theory is the Born approximation for scattering amplitudes. The problem itself is formulated as the scattering of a given particle on a spherically symmetric potential $V$, treating its wave function as infinite, thereby reducing the problem to the stationary case. At large distances we obtain incoming and scattered plane waves, described spatially by $f(\theta, \phi)$. This quantity is a probability amplitude which, on taking the squared modulus, gives the differential cross-section. In general the problem is solved by writing the Schrödinger equation in the implicit form known as the Lippmann--Schwinger equation, and then solving it iteratively. Restricting oneself to the first iteration is called the Born approximation and has the form:

\[f(\theta) = \frac{2m}{K\hbar}\int_0^\infty dr' r'V(r')\sin(Kr') \qquad \merr\]

where $K=\left|\vec k_0 - k_0\frac{\vec r}{r}\right|$.

\section{The Schrödinger equation as the basis for understanding the structure of the atom and of the chemical molecule; the periodic table of the elements, chemical bonds}

\subsection{The Schrödinger equation}

The Schrödinger equation is an equation of the form:

\begin{equation}
    i\hbar\frac{d}{dt}\ket{\psi\left(t\right)} = \hat{H}\ket{\psi\left(t\right)}
\end{equation}

% NOTE: the source has a corrupted word ("kota skladowa"); rendered here as the
% z-component, which is what is meant.
In the course of solving the equation it is found that it has meaning only for a specific set of natural numbers---the quantum numbers: the principal $n$, the azimuthal $l$ and the magnetic $m$. This is equivalent to showing that the energy of the electron, the square of the angular momentum and the $z$-component of the angular momentum are quantised. Each of the eigenfunctions so obtained, $\Psi_{nlm}\left(r,\theta,\phi\right)$, is an orbital.

In chemistry, the \textbf{electron shell} around a given atom is taken to be the set of atomic orbitals having the same principal quantum number $\mathsf{n}$.

The Schrödinger equation also permits the binding energy of the electron to be determined.

\subsection{The periodic table of the elements}

% NOTE: the source is internally inconsistent here --- the table is ordered by
% atomic number, while the periodic law is stated in terms of atomic masses
% (Mendeleev's original, pre-Moseley formulation).
\textbf{The periodic table of the elements} --- a compilation, in the form of a table, of all the chemical elements, ordered according to their increasing atomic number, grouping the elements according to their cyclically recurring similarities of properties, in accordance with \textbf{Dmitri Mendeleev's periodic law}. It states that the properties of the elements are periodically dependent on their atomic masses.\merr The vertical columns in the periodic table are groups. Chemical similarity within a group is explained by the identical or similar electron configuration of the valence shell of the atoms of the elements of that group. The number of the period (the horizontal row), in turn, corresponds to the number of the valence shell of the atom.

\section{Quantum tunnelling; examples}

Quantum tunnelling is the phenomenon of a particle passing through a potential barrier of height greater than the energy of the particle. It has no classical counterpart or analogy. In quantum mechanics a particle is not a rigid body of definite size and precisely specified position, although the interpretation of the wave function of a stationary state should incline one to the opposite conclusion (if the wave function is interpreted as a mechanical wave, then the particle should always be located in the middle of the potential well; the stationarity of the solutions at any instant would deprive the particle of any motion whatsoever, were this not accompanied by a change of energy levels). It is represented by a wave function which determines the probability of localising the particle in a specified region of space through the squared modulus of the wave function in that region.

The solution of the Schrödinger equation for a particle with a potential barrier allows the wave function to be found for a particle located outside the barrier. The wave function inside the potential barrier takes the form of waves travelling in two directions:

% NOTE: as in the source, both terms carry the same amplitude A; in general the
% forward- and backward-travelling amplitudes differ.
\begin{equation}
    \psi(x) = Ae^{i k_1 x} + Ae^{-i k_1 x} \qquad \merr
\end{equation}

In the middle of the potential barrier the wave function of the particle undergoes exponential damping, dependent on the height of the barrier and on the energy of the particle. Outside the potential barrier the wave function takes the form of a wave travelling outwards:

\begin{equation}
    \psi(x) = Be^{i k_2 x}
\end{equation}

where $k_1$ and $k_2$ are wave numbers dependent on the energy of the particle. Accordingly, the amplitude of the wave function on the far side of the barrier, $B$, is always smaller than the amplitude of the wave function inside, $A$. The tunnelling coefficient is expressed as $T = \frac{B^2}{A^2}$.

Nuclear fusion, the source of the Sun's energy, occurs to a large extent thanks to the tunnelling phenomenon. This phenomenon makes it possible to overcome the barrier of Coulomb repulsion of atomic nuclei at a temperature lower than would follow from the laws of thermodynamics. The tunnel effect also raises hopes of lowering the temperature of fusion carried out in a controlled manner. Thanks to the tunnelling phenomenon, the emission of $\alpha$ particles occurs in the radioactive decay of massive atomic nuclei.

In modern technology, the operation of many semiconductor electronic components (e.g.\ the tunnel diode) and of devices such as the scanning tunnelling microscope is based on the tunnelling phenomenon.

\section{Physical foundations of optical spectroscopy; electronic, vibrational and rotational excitations of molecules and the corresponding ranges of the electromagnetic spectrum}

Optical spectroscopy is an analytical technique which consists in studying the interaction of electromagnetic radiation with matter. By analysing the absorption, emission or scattering of light by atoms and molecules, information may be obtained about their electronic, vibrational and rotational structure.

\subsection{Electronic excitations}

Electronic excitations refer to transitions of electrons between bound energy levels in atoms or molecules. When an atom or molecule absorbs a photon of appropriate energy, an electron may jump from a lower energy level to a higher one, which is visible in the absorption spectrum. These transitions usually have frequencies in the range of visible light or the near ultraviolet.

\subsection{Vibrational excitations}

Vibrational excitations concern the vibrations of atoms in molecules. When a molecule absorbs electromagnetic radiation, its atoms may begin to oscillate about their equilibrium positions. In particular, the stretching and bending of chemical bonds may occur. The order of magnitude is smaller than in the case of electronic transitions and covers the infrared range.

\subsection{Rotational excitations}

Rotational excitations refer to the rotations of molecules about their centres of mass. When a molecule absorbs energy, its angular momentum may change, which causes a rotational excitation. Together with vibrational excitations they constitute an important source of knowledge about the structure of molecules, especially organic ones, although these spectra can be very complicated and harder to analyse. Rotational transitions occur in the far infrared and microwave range.

\section{Phenomena accompanying the interaction of electromagnetic radiation with atoms and molecules; stimulated and spontaneous emission, absorption}

\subsection{Absorption}

\textbf{Absorption} is the process of the taking up of the energy of an electromagnetic wave by a substance. The intensity of the light of a beam passing through the substance is thereby reduced.

In cold and rarefied \textbf{atomic gases} the energy levels characterising the excited states lie relatively far from one another. These substances may therefore absorb only certain photons of strictly defined energies. If the spectrum of the incident light is a continuous spectrum, this causes dark lines to appear in that spectrum.

In the case of \textbf{molecules} the system of electron energy levels is more complex, because besides the levels associated with the electron configuration there are additionally vibrational levels---associated with the vibrations of atoms within the molecule---and rotational levels---associated with rotations of the whole molecule. The energy levels lie so close to one another that they merge into whole bands spanning a given range of energy.

\subsection{Spontaneous emission}

\textbf{Spontaneous emission} --- the emission of photons occurring when electrons occupying excited levels return spontaneously to lower energy levels.

The phenomenon occurs universally and is responsible for almost every kind of luminescence of bodies, e.g.\ of heated gases, excited atoms, liquid and solid bodies, and also of electronic devices such as light-emitting diodes. In the absence of new excitations, the equation determining the number of atoms remaining in the excited state has the form:

\begin{equation}
    \mathsf{N}\left(t\right) = \mathsf{N}\left(0\right)e^{-\frac{t}{\tau}}
\end{equation}

where $\tau$ is the lifetime.

\subsection{Stimulated emission}

\textbf{Stimulated emission} --- the process of emission of photons by matter as a result of interaction with an initiating photon. The condition for stimulated emission to occur is the equality of the photon energy with the excitation energy of the atom. The photon initiating the emission is not absorbed by the matter---it plays only the role of triggering the process. The photon emitted by the atom has the same frequency (and hence also energy), phase and polarisation as the photon causing the emission. The direction of motion of both photons is also the same. Light composed of such identical photons is called coherent light. This phenomenon is the basis of the operation of lasers.

\section{Atoms in strong electric and magnetic fields (the Stark effect, the Zeeman effect); experimental methods of studying spin}

\paragraph{The Stark effect}

The Stark effect is the phenomenon consisting in the splitting of energy levels in an external electric field. Let us define the electric field along the $z$ axis:

\begin{equation}
    \vec{E} = E\hat{e}_z \xrightarrow[]{} V = -Ez = -Er\cos(\theta)
\end{equation}

This corresponds to the perturbing Hamiltonian $\hat{H}' = eEr\cos(\theta)$, so we may calculate the first-order correction to the ground state $\ket{nlm} = \ket{100}$:

\begin{equation}
    E^{(1)}_{100} = \bra{100}\hat{H}'\ket{100} = eE\int\psi^*_{100}r\cos(\theta)\psi_{100}d^3r = 0
\end{equation}

Since the ground state has $l=0$ it is spherically symmetric, and consequently the first-order correction $E^{(1)}_{100} = 0$.

% NOTE: the source calls what follows a second-order correction, but the
% procedure carried out is degenerate first-order perturbation theory within
% the n=2 manifold.
We may, however, calculate the correction in second-order perturbation theory\merr, i.e.\ find the matrix elements $\bra{2lm}\hat{H}'\ket{2lm}$. The only non-zero elements are:

\begin{equation}
    \bra{200}\hat{H}'\ket{210} = \bra{210}\hat{H}'\ket{200} = -3eEa_0
\end{equation}

\begin{equation}
\hat{H}' =
    \begin{pmatrix}
        0 & -3eEa_0 & 0 & 0 \\
        -3eEa_0 & 0 & 0 & 0\\
         0 & 0& 0 & 0\\
         0 & 0& 0 & 0\\
    \end{pmatrix}
\end{equation}

Calculating the eigenvectors and eigenvalues of this matrix we obtain $\lambda_1 = \lambda_2 = 0$, $\lambda_3 = -\lambda_4 = 3eEa_0$. In the presence of an electric field two states arise:

\begin{equation}
    \frac{1}{\sqrt{2}}(\ket{200}\pm\ket{210})
\end{equation}

with energies:

% NOTE: the source writes E^{(0)}_{100}; since the states involved belong to the
% n=2 manifold, this should be E^{(0)}_{200}.
\begin{equation}
    E = E^{(0)}_{200} \pm 3eEa_0
\end{equation}

\paragraph{The Zeeman effect}

The Zeeman effect is the phenomenon consisting in the splitting of energy levels in an external magnetic field. Most energy levels in atoms and molecules are degenerate with respect to energy---this means that there exist several levels of the same energy. What distinguishes the degenerate levels are other physical quantities characterising the electron in the atom---the orbital angular momentum, the spin, and so on. The presence of an external magnetic field may lift the degeneracy, because it interacts in different ways with electrons of different magnetic moment. A separation of the energy levels then takes place. In consequence, in the observed pattern of radiation emitted during electronic transitions, several closely spaced spectral lines appear in place of a single line.

Energy correction --- we introduce the Hamiltonian perturbing by a field $B$ in the $z$ direction:

\begin{equation}
    \hat{H}' = -\mu B_z
\end{equation}

where $\mu$ is the magnetic moment:

% NOTE: as in the source, the braces do not render as a bracket and L, S are
% vectors: \mu = -\frac{\mu_B}{\hbar}\left(\vec{L} + g_S\vec{S}\right).
\begin{equation}
    \mu=-\frac{\mu_B}{\hbar}\left(\vec{L} + g_S \vec{S}\right)
\end{equation}

$\mu_B$ is the Bohr magneton and $g_S \approx 2$ is the Landé factor for the electron. In the case where $S = 0$ we have the so-called normal Zeeman splitting, and the correction to the energy is:

% NOTE: the superscript should be (1), not (0); moreover, since the eigenvalue of
% L_z is \hbar m_l, the \hbar in the denominator cancels and the correction is
% -\mu_B B m_l. Retained as in the source.
\begin{equation}
    E^{(0)}_{nlm_l} = \bra{nlm_l}\hat{H}'\ket{nlm_l} = -\frac{\mu_B}{\hbar}Bm_l \qquad \merr
\end{equation}

In the case where $S\neq 0$ we have the anomalous Zeeman effect, and the correction to the energy is

\begin{equation}
    E^{(0)}_{nlm_lm_s} = -\frac{\mu_B}{\hbar}B(m_l + 2m_s) \qquad \merr
\end{equation}

\paragraph{Experimental methods of studying spin}

One of the best known methods of experimentally studying spin is the Stern--Gerlach experiment. In its original version, the experiment consisted in passing a beam of silver atoms through an inhomogeneous magnetic field and observing the image of the beam on a screen. Charged fundamental particles with non-zero spin, and also the bound states known to us with non-zero spin, have a non-zero magnetic moment whose magnitude and projection onto a chosen axis are related to the magnitude and projection of the spin. An unpolarised beam of such particles may be passed through an inhomogeneous magnetic field, and its splitting into several beams observed, depending on the value of the total spin. According to quantum mechanics, in the experiment with silver atoms we should therefore register atoms reaching only two points of the screen, in the case of a perfectly homogeneous beam of atoms with very small transverse dimensions. In practice, on account of the finite transverse dimensions of the beam and the unavoidable spread of atomic velocities, two separated spots are obtained.

\section{Thermal radiation; the emissive and absorptive power of bodies, Kirchhoff's law, black-body radiation}

Thermal radiation is a kind of electromagnetic radiation emitted by all bodies at a temperature above absolute zero. It is the result of the thermal motion of particles in the material and spans a wide range of wavelengths, from the infrared to visible radiation and even the ultraviolet. The basic concept is that of the black body, which is an ideal emitter and absorber of thermal radiation. It absorbs all radiation incident upon it regardless of wavelength, direction and polarisation, and emits radiation with the maximum possible efficiency for a given temperature.

\subsection{Emissive and absorptive power}

\textbf{Emissive power} ($E(\nu, T)$) determines the capacity of a body to emit radiation and is equal to the energy of the radiation emitted at a given frequency and temperature per unit area.

\textbf{Absorptive power} ($A(\nu,T)$) is the ratio of the power absorbed by the body to the incident power. Together with the reflection and transmission coefficients it sums to unity.

\subsection{Kirchhoff's law}

Kirchhoff's law of thermal radiation states that at a fixed temperature the function $\varepsilon(\nu,T)=\frac{E(\nu,T)}{A(\nu,T)}$ is universal for all bodies. Since for a black body $A=1$, this function is the emissive power of the black body.

\subsection{The black body}

The formula for the emissivity of a black body is the Planck distribution, which arose initially from fitting a function and was later derived on the assumption of the quantum nature of photons, and is given by:

\[\varepsilon(\nu,T) = \frac{2\pi\nu^2}{c^2}\frac{h\nu}{e^{h\nu/k_BT} - 1}\]

Auxiliary laws describing black-body radiation are: Wien's law, giving the wavelength at which it attains its maximum as depending on the temperature according to $\lambda_{max} = b / T$, and the Stefan--Boltzmann law, giving the total power emitted at a given temperature, $\Phi = \sigma T^4$.

\section{Diffraction and interference of waves; coherence, diffraction of radiation on crystals}

\subsection{Interference and coherence}

The phenomenon of interference, that is, of the superposition of waves, occurs only in the case where the interfering waves are coherent.

\textbf{Coherence} --- the difference of the phases arriving at a given point is constant in time, e.g.\ a laser.

The interference of waves emitted by two coherent sources was demonstrated in Young's experiment. A beam of light originating from a single source, having passed through one narrow slit, falls on two further slits separated from one another by $d$. We treat these slits as two point sources of coherent spherical waves. On a screen located at a large distance from the slits, dark and bright fringes are recorded. The waves originating from both slits are described by the wave functions:

% NOTE: the second line in the source is also labelled \psi_1; it should be \psi_2.
\begin{equation}
    \begin{matrix}
        \psi_{1}\left(r_{1},t\right) = A\left(r_{1}\right)\cos{\left(kr_{1} - \omega t\right)}\\
        \psi_{2}\left(r_{2},t\right) = A\left(r_{2}\right)\cos{\left(kr_{2} - \omega t\right)}
    \end{matrix}
\end{equation}

On the screen we observe the result of the summation of the waves:

\begin{equation}
    \psi_{1}\left(r_{1},t\right) + \psi_{2}\left(r_{2},t\right) = 2A\left(r_{av}\right)\cos{\left(k\frac{r_{1} - r_{2}}{2}\right)}\cos{\left(kr_{av} - \omega t\right)}
\end{equation}

The first factor (the cosine) describes the spatial modulation, while the second term, with the amplitude, describes the spherical wave. On the screen we obtain interference maxima (bright fringes) when the following condition is satisfied:

\begin{equation}
    \begin{matrix}
        \cos{\left(k\frac{r_{1} - r_{2}}{2}\right)} = \pm 1\\
        k\frac{r_{1} - r_{2}}{2} = m\pi
    \end{matrix}
\end{equation}

That is, the difference of the optical paths of the two rays is equal to a multiple of the wavelength: $r_{1} - r_{2} = m\lambda$, where $m$ is an integer. Cancellation of the waves (dark fringes) occurs when:

\begin{equation}
    \begin{matrix}
        \cos{\left(k\frac{r_{1} - r_{2}}{2}\right)} = 0\\
        k\frac{r_{1} - r_{2}}{2} = \frac{2m+1}{2}\pi
    \end{matrix}
\end{equation}

and thus for a difference of the optical paths of both rays equal to an odd multiple of half the wavelength: $r_{1} - r_{2} = \frac{2m+1}{2}\lambda$.

\subsection{Diffraction}

There is no essential difference between the phenomena of diffraction and interference. It has been agreed that the changes of intensity arising as a result of the superposition of waves produced by a finite number of discrete coherent sources are called interference, while those produced by coherent sources distributed in a continuous manner are conventionally called diffraction. In general one may say that diffraction is the set of phenomena arising during the propagation of waves in a medium with sharp inhomogeneities---if a wave encounters an obstacle on its path, then that part of the wave which passes beyond the obstacle will become a source of new spherical waves (Huygens' principle), that is, will undergo bending (diffraction). At every point beyond the obstacle there occurs a superposition of these partial waves, taking into account the amplitudes and phases of the waves.

\subsection{Diffraction of radiation on crystals}

\textbf{Bragg's law} --- a relation linking the structure of a crystal with the wavelength of the radiation incident on the crystal and the angle at which the interference maximum is observed. This law concerns so-called Bragg diffraction. When X-radiation falls on a crystal, diffraction of this radiation occurs at each of its atoms. The Bragg condition assumes diffraction on the planes on which the atoms of the crystal are arranged. With a known interplanar spacing and wavelength, Bragg's law determines the angle at which the wave must fall on the crystal for constructive interference (reinforcement) to occur. This means that for X-rays falling on a crystal the maxima of the diffracted radiation occur only for certain angles of incidence.

\begin{equation}
    n\lambda = 2d\sin{\theta}
\end{equation}

where $n$ is the order of diffraction, $\lambda$ the wavelength of the radiation, $d$ the interplanar spacing---the distance between the planes on which the scattering occurs---and $\theta$ the glancing angle---the angle of incidence defined as the angle between the beam of primary rays and the plane of the crystal.

A crystal consists of many parallel planes on which atoms are arranged---in an identical manner on each of the planes---and each pair of neighbouring planes is separated from one another by a constant distance. The beam under consideration falls on the crystal and is diffracted on each of the planes, and as a result many scattered beams appear (there are as many of them as there are planes), but all of them are scattered in the same direction, and their interference takes place. The diffraction of the incident beam occurs at each of the atoms, but since these are arranged periodically in planes, one considers the diffraction on the planes. The wave diffracted on a single atom is spherical, but for atoms lying on a given plane the fronts of these spherical waves form the front of a plane wave, and thus the beam scattered---diffracted---on that whole plane. Let us consider the part of the incident beam falling on a given plane and diffracted on that plane, and the part of the incident beam falling on another plane and diffracted on that other plane. If the difference of the optical paths for such pairs of incident and reflected beams (equal in this case to the difference of the geometrical paths) is equal to an integer multiple of the wavelength, then as a result of interference a reinforcement of the diffracted beam will occur---and it is precisely this case that the Bragg condition describes. On account of the fact that the angle at which the reinforcement occurs is equal to the angle of incidence of the radiation, the phenomenon discussed is often called reflection---however, diffraction on a crystal differs from reflection:

\begin{itemize}
    \item the beam diffracted on a crystal consists of rays scattered on all the atoms of the crystal, whereas in ordinary reflection this occurs only for the surface layer;
    \item diffraction on a crystal occurs only for particular angles of incidence determined by the Bragg equation, whereas reflection occurs for all angles;
    \item the intensity of the beam diffracted on a crystal is much smaller than the intensity of the incident beam, whereas good mirrors reflect close to 100 per cent of the light falling on them.
\end{itemize}

\section{Indistinguishability of elementary particles; the connection between spin and statistics, Bose--Einstein condensation}

In the case of identical elementary particles, e.g.\ electrons, quarks or photons, the wave function of any system of them is written in the form:

\begin{equation}
    \psi (1,2,\dots,N,t)
\end{equation}

where $(1,2,\dots,N)$ denotes the set of position and spin coordinates of the individual particles of the quantum system, the remaining parameters having to be the same. If it is assumed that these particles have been labelled unambiguously, so that each of them has an individual number, then on exchanging two particles (e.g.\ the first and the second) the wave function may change only by a phase factor $e^{i\varphi}$, i.e.:

\begin{equation} \label{eq:identical}
    \psi (1,2,\dots,N,t) = e^{i\varphi} \psi (2,1,\dots,N,t)
\end{equation}

where $\varphi$ is an arbitrary real number, and such a change does not cause a change in the modulus of the wave function. When the particles are exchanged back, the phase factor changes again and we obtain

\begin{equation}
    \Tilde{\psi} (2,1,\dots,N,t) = e^{i\varphi} \psi (1,2,\dots,N,t)
\end{equation}

Returning to equation \ref{eq:identical} we obtain:

\begin{equation}
    \Tilde{\psi} (1,2,\dots,N,t) = e^{i\varphi}e^{i\varphi} \psi (1,2,\dots,N,t)
\end{equation}

Since the wave functions $\Tilde{\psi}$ and $\psi$ correspond to the same configuration of particles, their moduli must be equal to one another, and therefore $e^{2i\varphi} = \pm 1$. Consequently, the wave function of identical particles must be either symmetric or antisymmetric.

\paragraph{The spin--statistics postulate}

The wave function of a system of indistinguishable particles with half-integer spin (i.e.\ 1/2, 3/2, etc.) must be antisymmetric with respect to the exchange of the spatial and spin coordinates of any pair of particles. Analogously, for particles of integer spin (i.e.\ 0, 1, 2, etc.) it must be symmetric. It follows from this that two fermions (particles of half-integer spin) cannot be in the same quantum state (the Pauli exclusion principle).

\paragraph{Bose--Einstein condensation}

The statistics refers to a non-interacting gas of bosons at low temperatures. Bosons are not subject to the Pauli exclusion principle, and therefore many particles may occupy the same state within a single system. On the basis of these assumptions one obtains the formula for the occupation density of the $k$-th state:

\begin{equation}
    n_k = \frac{1}{e^{(\varepsilon_k - \mu)/kT} - 1}
\end{equation}

$\varepsilon_k$ denotes the energy of the $k$-th state, $\mu$ the chemical potential, $k$ the Boltzmann constant and $T$ the temperature. It is worth emphasising that the index $k$ corresponds to the wave vector $\vec{k}$. On account of the difference appearing in the denominator, under appropriate conditions the density of the lowest states is very large. The effect of condensation is the collective behaviour of all the particles taking part in it (approximately, they all behave like a single particle). This is not a matter of condensation in the ordinary sense in position space---the particles are not located in one place---but of the ``condensation'' of particles in momentum space: a considerable number of particles have the same momentum. The critical temperature is given by the formula:

\begin{equation}
    T_c \approx 3.3125 \frac{\hbar^2 n^{2/3}}{m k_B}
\end{equation}

where $n$ is the density of atoms and $m$ the mass of a single atom.

\section{Statistical ensembles; classical and quantum statistics, thermodynamic potentials and their connection with the partition functions of the corresponding ensembles}

Statistical ensembles are a concept used in statistical mechanics to describe the macroscopic properties of physical systems on the basis of the statistics of the microscopic states of the particles. Statistical ensembles are hypothetical collections of very many copies of the same system, where each copy may be in a different microscopic state consistent with the adopted macroscopic conditions. To begin with, we shall assume that the fixed macroscopic parameters will be the number of particles $N$, the energy $E$ and the volume $V$.

\subsection{The microcanonical ensemble}

All microstates are equally probable, and their partition function $\Omega$ is simply the number of all microstates. Boltzmann connected such a sum with the entropy by means of the formula

\[S=k_B \log\Omega\]

\subsection{The canonical ensemble}

% NOTE: the source writes the sum as \sum_{r=1}^{N}; the sum runs over all
% microstates r, not over the N particles.
Imposing a constraint on the energy significantly complicates the calculations, which may be simplified by introducing a constraint on the temperature and assuming that the system exchanges energy with a heat reservoir at constant temperature $T$. The microstates are then subject to the Boltzmann distribution, and the partition function is given by

\[Q_N(V,T) = \sum_{r}\exp\left(-E_r / k_BT\right)\]

and gives directly the Helmholtz free energy

\[A=-k_BT\log Q_N\]

\subsection{The grand canonical ensemble}

Just as summing over all energies simplified the calculations relative to the microcanonical ensemble, so here we are dealing with a further change of constraint. The number of particles is also difficult to define or measure, and so here too it becomes a variable, while the chemical potential $\mu$ is constrained. Counting over all possible numbers of particles with the generalised Boltzmann distribution, we obtain the partition function

\[Q=\sum_{r,s}\exp\left(-\alpha N_r - \beta E_s\right) \]

where $\beta = 1/k_BT$ and $\mu = - \alpha/\beta$. This sum is connected with the grand thermodynamic potential by:

\[ \Xi = -k_BT\log Q\]

\subsection{Quantum statistics}

The theory of statistical ensembles may be successfully applied also within the formalism of quantum mechanics. Considering the case of a system of $N$ free particles, and hence by means of the canonical ensemble, additional constraints are however imposed on the occupations of the energy states. This leads to the formulation of three statistics:

\begin{enumerate}
    \item the classical one, where the occupations are given by the Maxwell--Boltzmann distribution. This is an approach analogous to the above, but it does not explain the phenomena observed for a gas of electrons (fermions) or photons (bosons). This distribution does not assume that the particles are indistinguishable, but on account of the scale this has no significant effect on the correctness of the results.
    \item Fermi--Dirac, which takes into account the Pauli exclusion principle, to which fermions are subject. It states that a given state may be occupied by only one particle. This gives the following occupation distribution:
    \[\langle n_\varepsilon\rangle=\frac{1}{e^{\beta(\varepsilon - \mu)} + 1}\]
    \item Bose--Einstein, which describes particles not subject to the Pauli exclusion principle, but still indistinguishable. The distribution is given by:
    \[\langle n_\varepsilon\rangle=\frac{1}{e^{\beta(\varepsilon - \mu)} - 1}\]
\end{enumerate}

In the classical limit both quantum statistics converge to the Maxwell--Boltzmann distribution.

\section{Statistical explanation of the macroscopic properties of systems of very many particles: extensive and intensive parameters of the state of a system, the equation of state, specific heats (of gases and solids), magnetic susceptibilities, etc.}

\subsection{Extensive and intensive parameters}

The parameters used in the description of thermodynamic systems may be:

\begin{itemize}
    \item intensive, i.e.\ independent of the size of the system, e.g.\ pressure, temperature, concentrations of the components, molar volumes, molar internal energies, entropies and other molar quantities
    \item extensive, i.e.\ dependent on the size of the system, e.g.\ volume, masses or numbers of moles of the components, total volume, internal energy, entropy
\end{itemize}

\subsection{The equation of state}

An equation of state is a relation between the parameters (state functions) of a thermodynamic system, such as:

\begin{itemize}
    \item pressure
    \item density
    \item temperature
    \item entropy
    \item internal energy
\end{itemize}

which may be written in the form:

\begin{equation}
    p = p\left(\rho,T,s,u,...\right)
\end{equation}

The most commonly used equations of state in thermodynamics are the ideal gas equation of state and the van der Waals equation:

% NOTE: two issues carried over from the source: the ideal gas law is written
% with a lower-case v (should be pV = nRT), and the van der Waals equation mixes
% the molar form (a/V^2, V-b) with nRT on the right-hand side. The consistent
% n-mole form is (p + an^2/V^2)(V - nb) = nRT.
\begin{equation}
    \begin{matrix}
        \mathsf{pV} = \mathsf{nRT}\\
        \left(\mathsf{p} + \frac{\mathsf{a}}{\mathsf{V}^{2}}\right)\left(\mathsf{V}-\mathsf{b}\right) = \mathsf{nRT} \qquad \merr
    \end{matrix}
\end{equation}

where the parameter $b$ accounts for the volume of the molecules, while the parameter $\frac{a}{V^{2}}$ accounts for their mutual interaction.

\subsection{Specific heats}

\textbf{Specific heat} --- the heat needed to change the temperature of a body of unit mass by one unit

\begin{equation}
    \mathsf{c} = \frac{\Delta \mathsf{Q}}{m\Delta T}
\end{equation}

The heat capacity of a gas at constant parameter $\mathsf{x}$ may be calculated from the formula:

\begin{equation}
    c_{x} = \left(\frac{\partial \mathsf{Q}}{\partial\mathsf{T}}\right)_{x}
\end{equation}

The most commonly encountered heat capacities are those at constant volume and at constant pressure, $c_{V}$ and $c_{p}$. For a monatomic ideal gas:

\begin{equation}
    \begin{matrix}
        c_{V} = \frac{3}{2}R\\
        c_{p} = \frac{5}{2}R
    \end{matrix}
\end{equation}

\subsection{Magnetic susceptibility}

For magnetic systems we have the relation:

\begin{equation}
    \mathsf{dW} = \mu_{0}\mathsf{H}\mathsf{dM}
\end{equation}

where $\mathsf{M}$ denotes the magnetic moment. Magnetisation is the magnetic moment per unit volume:

\begin{equation}
    \mathsf{m} = \frac{\mathsf{M}}{V}
\end{equation}

The magnetic susceptibilities at constant temperature and the adiabatic one may be calculated:

\begin{equation}
    \begin{matrix}
        \chi_{T} = \frac{1}{\mu_{0}}\left(\frac{\partial\mathsf{M}}{\partial\mathsf{H}}\right)_{T}\\
        \chi_{S} = \frac{1}{\mu_{0}}\left(\frac{\partial\mathsf{M}}{\partial\mathsf{H}}\right)_{S}
    \end{matrix}
\end{equation}

\section{The Dirac and Klein--Gordon equations; the form of the equations and examples of objects governed by each of them}

The Klein--Gordon equation may be derived by using the representation of energy and momentum as operators: $E \xrightarrow[]{} i\hbar \frac{\partial}{\partial t}$ and $\vec{p}\xrightarrow[]{} - i\hbar\nabla$. Substituting the operators so obtained into the relativistic energy relation:

\begin{equation}
    E^2 = p^2c^2 + m^2c^4
\end{equation}

we obtain:

\begin{equation}
    \left( \Box + \frac{m^2c^2}{\hbar^2}  \right)\psi(\vec{r},t) = 0
\end{equation}

This equation describes scalar or pseudoscalar particles of zero spin. The simplest solution of the Klein--Gordon equation is a plane wave $\psi \propto e^{ik^jx^j - i\omega_k t}$, giving the relativistic dependence of the energy $\varepsilon_k = \hbar\omega_k$ on the momentum in the form $\varepsilon_k = \pm c\sqrt{p^2 + m_0^2c^2}$. The negative-energy solution of the Klein--Gordon equation has no direct physical meaning. This is caused by the erroneous assumption that relativistic wave equations can describe the dynamics of relativistic particles.

% NOTE: this "mass increase" formula could not be verified against a standard
% source; treat the numerical form with caution.
One of the applications of the Klein--Gordon equation is the determination of the increase of the mass of an electron in the field of an electromagnetic wave $A$ of field energy density $\rho$.

\begin{equation}
    M^2 = m_0^2 + \frac{\hbar^3e^2\rho}{\omega c^2} \qquad \merr
\end{equation}

It is worth noting that when the mass $m_0$ is formally equal to zero, i.e.\ the Klein--Gordon equation describes a massless particle, this particle acquires a mass through the interaction with the electromagnetic field $A$ alone.

\paragraph{The Dirac equation}

The Dirac equation for a free particle takes the form:

\begin{equation}
(i\hbar \gamma^{\mu}\partial_{\mu} - mc)\Psi(x^{\nu}) = 0
\end{equation}

where $\gamma^{\mu}$ are the Dirac gamma matrices and $\partial_{\mu}$ are the components of the four-gradient. The Dirac matrices must satisfy the relation $\{\gamma^{\mu},\gamma^{\nu}\} = 2g^{\mu \nu }I$, where the most common representation is: % NOTE: in the source the anticommutator braces were not escaped and did not render.

\begin{align}
    \gamma^0 =
    \begin{pmatrix}
        I & 0\\
        0 & -I
    \end{pmatrix}\\
    \gamma^i =
    \begin{pmatrix}
        0 & \sigma_i\\
        -\sigma_i & 0
    \end{pmatrix}
\end{align}

where $\sigma_i$ are the Pauli matrices. The equation is valid for particles of arbitrarily large energies (so-called relativistic particles) with spin 1/2 (fermions, e.g.\ electrons, quarks), both free and interacting with the electromagnetic field. The existence of spin follows from the very requirement of relativistic invariance of the equation of motion of particles. The Dirac equation was confirmed in relation to the fine structure of the spectrum of the hydrogen atom, showing excellent agreement with measurements. It predicts the existence of antiparticles.

\section{The structure of atomic nuclei and nuclear reactions; nuclear binding energy, fission and fusion reactions, alpha and beta decay, radioactive decay series, the limits of the world of nuclides}

The atomic nucleus consists of protons and neutrons, which are called nucleons. The number of protons in the nucleus defines the chemical element, and the number of neutrons determines the isotopes of that element. The nucleons are bound in the nucleus by nuclear forces, which are considerably stronger than the electromagnetic forces but act over a very short range (about 1--2 femtometres).

The binding energy is the energy needed to separate an atomic nucleus into its constituents (protons and neutrons). It is a measure of the stability of the nucleus. The greater the binding energy, the more stable the nucleus is, since it requires more energy to break it up into individual nucleons.

There exist various models of the structure of the atomic nucleus. For a long time the liquid drop model was dominant. It assumes that the nucleons in the nucleus behave like molecules in a drop of liquid, and consequently that the properties of the nucleus as a whole should be similar to the properties of a liquid drop. The microscopic interactions---the strong nuclear interaction and the electrostatic forces---are represented in this model by analogy with the forces of viscosity and surface tension. The formulae obtained in this way predict a constant binding energy per nucleon for light nuclei and a smaller one for nuclei of large mass. This leads to the conclusion that in large nuclei a separation into two fragments may occur, which explains the phenomena of fission of the atomic nuclei of heavy elements. This model is very approximate and does not explain all the properties of nuclei.

It was replaced by the shell model of the atomic nucleus, which arose by analogy with the shell model of the atom and, in accordance with observations of the excitation levels of atomic nuclei, assumes that nucleons within the nucleus cannot take arbitrary energy states, but only those consistent with the energies of successive shells. The model explains the deviations of the binding energies of nuclei from the energy determined in the liquid drop model. It also explains the existence of the ``magic numbers'': 2, 8, 20, 28, 50, 82, 126, for which atomic nuclei are the most stable. If a nucleus has one nucleon fewer or more, then its binding energy is markedly smaller.

Nuclear reactions are processes in which transformations of atomic nuclei occur under bombardment with particles (e.g.\ neutrons, protons):

\begin{itemize}
    \item fission --- a process in which a heavy nucleus divides into two (or more) lighter nuclei, emitting neutrons and releasing a large amount of energy.
    \item fusion --- two light nuclei combine to form a heavier nucleus, releasing energy.
    \item alpha decay --- the emission of an alpha particle (a helium-4 nucleus) by a nucleus.
    \item beta decay --- the emission of a beta particle (an electron or a positron) and an antineutrino/neutrino, as a result of the weak interaction.
    \item gamma decay --- the emission of high-energy electromagnetic radiation (gamma photons) by a nucleus. It usually follows alpha or beta decay, when the nucleus is in an excited state and returns to the ground state.
\end{itemize}

% NOTE: the neptunium series is usually classed as extinct in nature (its longest
% half-life, ~2.14 Myr, is short on geological timescales); the source calls all
% four "natural".
Radioactive decay series are sequences of successive radioactive decays which lead from one isotope to another until a stable isotope is reached. There are four natural radioactive series: the uranium, thorium, actinium and neptunium series.\merr

The limits of the world of nuclides are set by the proton and neutron drip lines and by the stability of superheavy nuclei.

% NOTE: the two sentences in the first item are in tension with one another
% (bismuth is Z = 83, yet the limit is said to lie beyond reach for Z > 30);
% retained as in the source.
\begin{itemize}
    \item \textbf{The proton drip line}: experimentally the proton drip line has been reached as far as bismuth (Z = 83). However, thanks to the Coulomb barrier, many nuclides beyond this line have long lifetimes and undergo beta transitions. The limit of stability with respect to proton emission is still far beyond the reach of experiments for $Z > 30$.\merr
    \item \textbf{The neutron drip line}: it has been reached almost up to N = 32. For larger N it is far beyond the reach of contemporary experiments. The structure of nuclei with a great excess of neutrons conceals many puzzles and remains the subject of intensive research.
    \item \textbf{The limit of stability for superheavy nuclei}: the limit for superheavy nuclei is spontaneous fission. A key role in their stability is played by shell effects, that is, by the microscopic structure of the nuclei. The existence of an island of superheavy, metastable nuclides is still an open question and constitutes one of the main challenges in research on nuclear matter.
\end{itemize}

\section{Nucleosynthesis and the energy sources of stars}

Three ways of the natural production of the elements may be distinguished (leaving decays aside). These are primordial nucleosynthesis, nucleosynthesis through thermonuclear reactions occurring in stars, and nucleosynthesis taking place in supernova explosions.

The first hydrogen nuclei (protons) came into being when the temperature had cooled sufficiently for quarks to be able to bind into hadrons. The period of hadron dominance is taken to be the period $10^{-6}$~s--1~s after the Big Bang (pBB). The further formation of the elements took place during primordial nucleosynthesis, which is dated to 10~s--$10^{3}$~s pBB. There arose then the nuclei of deuterium $(^{2}\mathrm{H})$, tritium $(^{3}\mathrm{H})$, helium ($^{3}\mathrm{He}$ and $^{4}\mathrm{He}$), lithium $(^{7}\mathrm{Li})$ and beryllium $(^{7}\mathrm{Be})$. The samples of beryllium-7 that are found are not regarded as originating from that period on account of its short half-life (53 days). In the course of decay, $^{7}\mathrm{Be}$ transformed into $^{7}\mathrm{Li}$.

Stars, which formed considerably later, were responsible for the next round of nucleosynthesis. Thanks to them the elements up to and including iron came into being. The primary substrate was hydrogen, which was burned into helium. Further, the helium nuclei were synthesised into atoms of heavier elements, up to iron.

It is often said that synthesis ends at iron, because the iron nucleus $^{56}\mathrm{Fe}$ is the most stable nucleus in nature. In reality $^{56}\mathrm{Fe}$ possesses the smallest mass per nucleon, while the most stable nucleus in nature is nickel $^{62}\mathrm{Ni}$. It possesses the greatest binding energy per nucleon (mass per nucleon $\neq$ binding energy per nucleon, because the mass of the proton $\neq$ the mass of the neutron). Synthesis processes in stars favour the production of the isotope $^{56}\mathrm{Ni}$ (instead of $^{62}\mathrm{Ni}$), which is less stable and decays to iron $^{56}\mathrm{Fe}$.

The largest stars, when they end their lives, explode as supernovae. In this process a great pressure arises and a great deal of energy is released. It is precisely then that the remainder of the periodic table known to us comes into being.

\subsection{Thermonuclear reaction}

A \textbf{thermonuclear reaction} is a phenomenon consisting in the joining of two lighter nuclei into one heavier nucleus. As a result of the exothermic reaction the energy released (in the form of the kinetic energy of the products and gamma radiation) is dissipated on the surrounding atoms and is transformed into thermal energy. The energy released during the reaction may be determined without carrying out the reaction, on the basis of the mass defect, that is, the difference of the masses of the reactants and the products of the reaction.

In not overly massive main-sequence stars the basic reaction is the synthesis of the helium nucleus. For synthesis to occur, the hydrogen nuclei (protons) must approach within the range of the nuclear interaction (about $1\,\mathrm{fm} = 10^{-15}$~m). The protons, however, repel one another electrostatically, and must therefore overcome a potential barrier of value about $E = 1$~MeV. Such thermal energy is possessed by particles at a temperature of $10^{10}$~K. Such a high temperature does not exist inside stars, but the course of the phenomenon at a lower temperature is explained by the tunnelling phenomenon.

\subsection{The proton--proton chain}

The \textbf{proton--proton chain} is a cycle of nuclear reactions in which a stable helium nucleus is formed from four hydrogen nuclei. During the transformations nuclear energy is released, which is the main source of the energy of the Sun and of other stars. The pp chain takes place in the cores of stars at temperatures from a few to a few tens of millions of kelvin.

\section{Classification of elementary particles; quarks and leptons and the particles transmitting the interactions, the structure of hadrons}

% NOTE: "quarks form hadrons in threes" contradicts the description of mesons
% (quark--antiquark) further down; retained as in the source.
The elementary particles building up matter are leptons and quarks. Leptons may occur on their own, whereas quarks form hadrons in threes.\merr The leptons are a group of 12 particles:

\begin{itemize}
    \item the electron and the electron neutrino
    \item the muon and the muon neutrino
    \item the tau and the tau neutrino
\end{itemize}

% NOTE: the source reads "neutrons" (neutronami) where it must mean "neutrinos";
% moreover, the charge of -1 applies to the charged leptons, while the
% corresponding antileptons carry +1.
and the corresponding antiparticles. Apart from the neutrinos, all of them have an electric charge equal to ``$-1$'' and interact weakly.\merr

There are likewise 12 quarks:

\begin{itemize}
    \item up (u) --- charge $+\frac{2}{3}$
    \item down (d) --- charge $-\frac{1}{3}$
    \item strange (s) --- charge $-\frac{1}{3}$
    \item charm (c) --- charge $+\frac{2}{3}$
    \item bottom (b) --- charge $-\frac{1}{3}$
    \item top (t) --- charge $+\frac{2}{3}$
\end{itemize}

and the corresponding antiparticles. Besides electromagnetic and weak interactions, they also interact strongly, which is connected with their colour charge (green, blue, red---these are only names and are in no way connected with optics). The bosons of this interaction are 8 massless gluons, which themselves also possess colour (they couple to themselves). It is necessary, in order to form a hadron, that the effective colour charge be neutral, which can be achieved by combining three quarks of three different colours (a baryon) or from two in pairs, e.g.\ red--antired (a meson). A characteristic quantum number is the baryon number (1 for baryons, 0 for mesons), which as yet is not associated with any symmetry following from deeper laws of physics (hence the searches for proton decay, the proton being the lightest hadron). Hadrons (baryons) are, for example, the proton and the neutron, together called ``nucleons''.

\section{Band structure of electronic states in crystals; metals, semiconductors, insulators}

A quantum-mechanical theory describing electrical conduction. In contrast to the classical theory, the starting point of this theory is Fermi--Dirac statistics and the wave nature of electrons. The most important concept of this theory is the energy band, that is, the range of energy that electrons in a conductor may possess. The existence of a continuous energy spectrum is connected with the interaction of the individual atoms with one another (it is a set of very closely spaced line spectra), while the occurrence of forbidden regions follows from the conditions imposed on the periodicity of the wave function of the electrons.

In an atom the individual electrons may be found in strictly defined, discrete energy states. Additionally, in a solid the atoms are bound to one another, which gives further restrictions on the admissible energies of the electrons. The allowed energy levels of isolated atoms are, as a result of interaction with other atoms in the crystal lattice, shifted, forming so-called allowed bands, i.e.\ the ranges of energy which electrons located in the individual orbits may take; levels lying outside the allowed bands are termed forbidden bands.

Electronics usually makes use of a simplified energy model, in which the energies of the valence electrons are described by two allowed bands:

\begin{itemize}
    \item the valence band --- the range of energy possessed by valence electrons bound to the atomic nucleus
    \item the conduction band --- the range of energy possessed by valence electrons freed from the atom, being then free carriers in the solid
\end{itemize}

There is also a forbidden band, where electrons cannot reside. In metallic conductors there is no forbidden band, which follows from two reasons:

\begin{itemize}
    \item The valence band is only partially filled with electrons and they can move freely, and so the valence band in conductors plays a role analogous to that of the conduction band in semiconductors and insulators
    \item The conduction and valence bands overlap, and therefore in the common band there are many free electrons and the flow of current is possible.
\end{itemize}

% NOTE: "GaAS" in the source is a typo for GaAs (band gap ~1.42 eV). The 0-6 eV
% and >6 eV thresholds are a convention of the source; other texts commonly place
% the semiconductor/insulator boundary nearer 3-4 eV.
Semiconductors possess a forbidden band between the valence band and the conduction band in the range from 0 to 6 eV\merr{} (for example Ge 0.7 eV, Si 1.1 eV, GaAs 1.4 eV, etc.).

It is assumed that at the temperature of absolute zero there are no electrons in the conduction band, whereas at higher temperatures electron--hole pairs arise; the higher the temperature, the more such pairs there are. Semiconductors may be doped by adding admixtures of other atoms into the structure of the crystal. In this way n-type semiconductors (donor-doped, ``giving up an electron'') and p-type semiconductors (acceptor-doped, ``accepting an electron'') arise.

Insulators have a forbidden band with an energy gap greater than 6 eV.\merr